\documentclass[12pt,oneside]{book}

%============================================================
% PACKAGES
%============================================================

%--------------------
% Page + font
%--------------------
\usepackage[margin=1in]{geometry}
\usepackage[T1]{fontenc}
\usepackage[utf8]{inputenc}
\usepackage{lmodern}
\usepackage{microtype}

%--------------------
% Math
%--------------------
\usepackage{amsmath, amssymb, amsfonts, amsthm, mathtools}
\usepackage{bm}
\usepackage{bbm}
\usepackage{mathrsfs}
\usepackage{dsfont}

%--------------------
% Tables + figures
%--------------------
\usepackage{graphicx}
\usepackage{float}
\usepackage{booktabs}
\usepackage{array}
\usepackage{xltabular}
\usepackage{multirow}
\usepackage{caption}
\usepackage{subcaption}

%--------------------
% Lists
%--------------------
\usepackage{enumitem}
\setlist[itemize]{topsep=2pt,itemsep=2pt}
\setlist[enumerate]{topsep=2pt,itemsep=2pt}

%--------------------
% Colors + hyperlinks
%--------------------
\usepackage[dvipsnames]{xcolor}
\usepackage[
    colorlinks=true,
    linkcolor=MidnightBlue,
    citecolor=MidnightBlue,
    urlcolor=MidnightBlue
]{hyperref}
\usepackage[nameinlink,capitalise]{cleveref}

%--------------------
% Algorithms
%--------------------
\usepackage{algorithm}
\usepackage{algpseudocode}

%--------------------
% Code blocks
%--------------------
\usepackage{listings}

\lstdefinestyle{codestyle}{
    basicstyle=\ttfamily\small,
    breaklines=true,
    frame=single,
    numbers=left,
    numberstyle=\tiny,
    keywordstyle=\color{MidnightBlue},
    commentstyle=\color{ForestGreen},
    stringstyle=\color{BrickRed},
    showstringspaces=false,
    tabsize=4
}

\lstset{style=codestyle}

%--------------------
% Section styling
%--------------------
\usepackage{titlesec}
\usepackage{tocloft}

\setcounter{tocdepth}{2}
\setcounter{secnumdepth}{3}

%============================================================
% THEOREMS
%============================================================

\numberwithin{equation}{section}

\theoremstyle{definition}
\newtheorem{definition}{Definition}[section]
\newtheorem{example}{Example}[section]
\newtheorem{exercise}{Exercise}[section]
\newtheorem{remark}{Remark}[section]

\theoremstyle{plain}
\newtheorem{theorem}{Theorem}[section]
\newtheorem{lemma}{Lemma}[section]
\newtheorem{proposition}{Proposition}[section]
\newtheorem{corollary}{Corollary}[section]

\theoremstyle{remark}
\newtheorem*{note}{Note}

%============================================================
% CUSTOM COMMANDS
%============================================================

% Sets
\newcommand{\R}{\mathbb{R}}
\newcommand{\N}{\mathbb{N}}
\newcommand{\Z}{\mathbb{Z}}
\newcommand{\Q}{\mathbb{Q}}
\newcommand{\C}{\mathbb{C}}

% Probability
\newcommand{\E}{\mathbb{E}}
\newcommand{\Var}{\mathrm{Var}}
\newcommand{\Cov}{\mathrm{Cov}}
\newcommand{\Corr}{\mathrm{Corr}}
\newcommand{\Prob}{\mathbb{P}}

\newcommand{\ind}{\mathbbm{1}}
\newcommand{\iid}{\overset{\text{iid}}{\sim}}
\newcommand{\convd}{\overset{d}{\to}}
\newcommand{\convp}{\overset{p}{\to}}
\newcommand{\as}{\overset{a.s.}{\to}}

% Operators
\newcommand{\given}{\,\middle|\,}
\newcommand{\argmin}{\operatorname*{arg\,min}}
\newcommand{\argmax}{\operatorname*{arg\,max}}

% Linear algebra / analysis
\newcommand{\norm}[1]{\left\lVert #1 \right\rVert}
\newcommand{\abs}[1]{\left\lvert #1 \right\rvert}
\newcommand{\inner}[2]{\left\langle #1, #2 \right\rangle}

\newcommand{\dd}{\,\mathrm{d}}
\newcommand{\grad}{\nabla}
\newcommand{\Hess}{\nabla^2}

\newcommand{\tr}{\operatorname{tr}}
\newcommand{\rank}{\operatorname{rank}}
\newcommand{\diag}{\operatorname{diag}}

% Distributions
\newcommand{\Normal}{\mathcal{N}}
\newcommand{\Bern}{\mathrm{Bernoulli}}
\newcommand{\Bin}{\mathrm{Binomial}}
\newcommand{\Pois}{\mathrm{Poisson}}
\newcommand{\Gam}{\mathrm{Gamma}}
\newcommand{\BetaDist}{\mathrm{Beta}}
\newcommand{\Exp}{\mathrm{Exponential}}
\newcommand{\Dir}{\mathrm{Dirichlet}}
\newcommand{\InvGam}{\mathrm{Inv\text{-}Gamma}}

% Title block
\title{EN.553.733 Nonparametric Bayesian Statistics}
\author{Jonathan Y. Ma}
\date{June 2026}

\begin{document}

\maketitle
\tableofcontents

\chapter{Introduction}
The Bayesian point of view with $y$ data and fixed $\theta$ paramaters:
\[
    \mathbb{P} (\theta | y) \propto \mathbb{P} (y|\theta)\mathbb{P} (\theta)
\]
Nonparametric Bayes (BNP) extends this by allowing unbounded/growing/infinite parameters

\section{Exchangeability and De Finetti's Theorem}

A theoretical motivation for nonparametric Bayesian methods comes from \emph{De Finetti's Theorem}.

\subsection*{Exchangeability}

A sequence $(X_n)_{n \ge 1}$ is \emph{infinitely exchangeable} if for every $N$ and every permutation $\sigma$ of $\{1,\ldots,N\}$,
\[
p(X_1, \ldots, X_N) = p\bigl(X_{\sigma(1)}, \ldots, X_{\sigma(N)}\bigr).
\]

This means the joint distribution depends only on the multiset of observations, not their order.

\subsection*{De Finetti Representation}

\begin{theorem}[De Finetti (informal)]
A sequence $X_1, X_2, \ldots$ is infinitely exchangeable if and only if, for all $N$, there exists a probability measure $P$ such that
\[
p(X_1, \ldots, X_N) = \int_{\Theta} \prod_{n=1}^N p(X_n \mid \theta)\, P(d\theta).
\]
\end{theorem}

Conditionally on $\theta$, the observations are i.i.d., i.e.,
\[
X_n \mid \theta \;\stackrel{i.i.d.}{\sim}\; p(\cdot \mid \theta).
\]

Thus exchangeability implies a mixture-of-i.i.d.\ representation.

\subsection*{Implications}

This representation motivates:
\begin{itemize}
\item likelihood models $p(x \mid \theta)$,
\item prior distributions $P(d\theta)$,
\item \emph{nonparametric Bayesian} priors, where $P$ is infinite-dimensional.
\end{itemize}

\section{The Bayesian Toolkit}

\subsection{Likelihood}

At the core of Bayesian inference is the \emph{likelihood}. Given data $D = x = (x_1,\ldots,x_n)$,
\[
L(\theta \mid D) = f(x \mid \theta).
\]

If $x_1,\ldots,x_n$ are conditionally independent given $\theta$, then
\[
f(x \mid \theta) = \prod_{i=1}^n f(x_i \mid \theta).
\]

\subsection{Prior Distribution}

Uncertainty about $\theta$ is encoded through a \emph{prior} $\pi(\theta)$.

\subsection{Joint, Marginal, and Posterior}

\paragraph{Joint distribution}
\[
\psi(\theta, x) = f(x \mid \theta)\pi(\theta).
\]

\paragraph{Marginal distribution}
\[
m(x) = \int f(x \mid \theta)\pi(\theta)\, d\theta.
\]

This quantity is also called the \emph{evidence} or \emph{model likelihood}.

\paragraph{Posterior distribution}
\[
\pi(\theta \mid x)
= \frac{f(x \mid \theta)\pi(\theta)}{\int f(x \mid \theta)\pi(\theta)\, d\theta}
= \frac{f(x \mid \theta)\pi(\theta)}{m(x)}.
\]

\subsection{Predictive Distribution}

If $y \sim g(y \mid \theta, x)$, then the predictive distribution is
\[
g(y \mid x) = \int g(y \mid \theta, x)\, \pi(\theta \mid x)\, d\theta.
\]

This integrates parameter uncertainty into predictions.

\paragraph{Remark}
While inference focuses on $\pi(\theta \mid x)$, prediction via $g(y \mid x)$ is often the primary goal.

\section{Conjugate Priors}

\subsection{Definition}

\begin{theorem}[Conjugate Prior]
A class $\mathcal{P}$ of prior distributions is conjugate for a likelihood $p(x \mid \theta)$ if
\[
\pi(\theta) \in \mathcal{P} \;\Rightarrow\; \pi(\theta \mid x) \in \mathcal{P}.
\]
\end{theorem}

\subsection*{Examples}

\begin{itemize}
\item If $x \mid \theta \sim \mathrm{Bin}(n,\theta)$ and $\theta \sim \mathrm{Beta}(a,b)$, then
\[
\theta \mid x \sim \mathrm{Beta}(a + x,\; b + n - x).
\]

\item If $x_1,\ldots,x_n \stackrel{i.i.d.}{\sim} \mathrm{Poisson}(\lambda)$ and $\lambda \sim \mathrm{Gamma}(a,b)$, then
\[
\lambda \mid x \sim \mathrm{Gamma}\bigl(a + \textstyle\sum_{i=1}^n x_i,\; b + n\bigr).
\]
\end{itemize}

\subsection{Properties}

\paragraph{Advantage}
\emph{Closed-form} posterior updates.

\paragraph{Disadvantage}
The conjugate family may not reflect prior beliefs.

\section{Credible Intervals}

\begin{theorem}[Credible Interval]
An interval $(a,b)$ is a $(1-\alpha)100\%$ credible interval for $\theta$ if
\[
\Pr(a < \theta < b \mid x) = 1 - \alpha.
\]
\end{theorem}

\subsection*{Quantile Construction}

Choose $a,b$ such that
\[
\Pr(\theta < a \mid x) = \frac{\alpha}{2},
\qquad
\Pr(\theta > b \mid x) = \frac{\alpha}{2}.
\]

\subsection*{Remark}

Credible intervals quantify posterior uncertainty, unlike frequentist confidence intervals which are defined via repeated sampling.

\section{MAP Estimator}

The \emph{maximum a posteriori (MAP)} estimator is defined as
\[
\hat{\theta}_{\mathrm{MAP}} = \arg\max_{\theta} f(x \mid \theta)\pi(\theta).
\]

Equivalently,
\[
\hat{\theta}_{\mathrm{MAP}} = \arg\max_{\theta} \left[ \log f(x \mid \theta) + \log \pi(\theta) \right].
\]

\subsection*{Motivation}

\begin{itemize}
\item Interpretable as a \emph{penalized likelihood} estimator.
\item Useful when optimization is easier than integration.
\end{itemize}

\section{Testing Hypotheses}

Hypothesis testing on $\theta$ is formulated as
\[
H_0: \theta \in \Theta_0
\qquad \text{versus} \qquad
H_1: \theta \notin \Theta_0.
\]

\section{Bayes Factor}

Bayesian testing procedures can be based on the posterior probability
\[
P^{\pi}(\theta \in \Theta_0 \mid x),
\]
or alternatively on the \emph{Bayes factor}.

\subsection*{Definition}

The Bayes factor comparing $\Theta_1$ to $\Theta_0$ is
\[
B^{\pi}_{10}
=
\frac{P^{\pi}(\theta \in \Theta_1 \mid x) / P^{\pi}(\theta \in \Theta_0 \mid x)}
{P^{\pi}(\theta \in \Theta_1) / P^{\pi}(\theta \in \Theta_0)}.
\]

\subsection*{Model Comparison Formulation}

For corresponding models $\mathcal{M}_1$ versus $\mathcal{M}_0$,
\[
B^{\pi}_{10}
=
\frac{P^{\pi}(\mathcal{M}_1 \mid x)}{P^{\pi}(\mathcal{M}_0 \mid x)}
\Big/
\frac{P^{\pi}(\mathcal{M}_1)}{P^{\pi}(\mathcal{M}_0)}.
\]

If the prior is written as
\[
\pi(\theta)
=
\Pr(\theta \in \Theta_1)\,\pi_1(\theta)
+
\Pr(\theta \in \Theta_0)\,\pi_0(\theta),
\]
then
\[
B^{\pi}_{10}
=
\frac{\int f(x \mid \theta_1)\pi_1(\theta_1)\, d\theta_1}
{\int f(x \mid \theta_0)\pi_0(\theta_0)\, d\theta_0}
=
\frac{m_1(x)}{m_0(x)}.
\]

Thus the Bayes factor is \emph{akin to a likelihood ratio}.

\section{Jeffreys' Scale}

Interpretation of the Bayes factor is often based on $\log_{10}(B^{\pi}_{10})$:

\begin{itemize}
\item if $\log_{10}(B^{\pi}_{10}) \in [0, 0.5]$, evidence against $H_0$ is poor,
\item if $\log_{10}(B^{\pi}_{10}) \in [0.5, 1]$, it is substantial,
\item if $\log_{10}(B^{\pi}_{10}) \in [1, 2]$, it is strong,
\item if $\log_{10}(B^{\pi}_{10}) > 2$, it is decisive.
\end{itemize}

\section{Markov Chain Monte Carlo (MCMC)}

Statistical inference is based on the posterior distribution
\[
\pi(\theta \mid x)
=
\frac{\pi(\theta) f(x \mid \theta)}{\int \pi(\theta) f(x \mid \theta)\, d\theta}.
\]

\subsection*{Motivation}

\begin{itemize}
\item In conjugate cases (e.g.\ normal-normal, binomial-beta), integrals can be computed analytically.
\item In general, closed-form expressions for $\pi(\theta \mid x)$ are not available.
\item We still need to compute:
\begin{itemize}
\item point estimates,
\item interval estimates,
\item hypothesis tests.
\end{itemize}
\end{itemize}

This motivates simulation-based methods.

\subsection*{The MCMC Principle}

\subsection*{Definition}

An \emph{MCMC method} for simulating from a distribution $\pi$ is any method that constructs an ergodic Markov chain $(X_n)$ whose stationary distribution is $\pi$.

Given an initial value $x^{(0)}$, the chain is generated via a transition kernel such that
\[
X_n \xrightarrow{d} X \sim \pi.
\]

\subsection*{Construction}

To build such a chain, common methods include:
\begin{itemize}
\item Metropolis--Hastings,
\item Gibbs sampling.
\end{itemize}

There are also diagnostic methods to assess whether the chain has reached its stationary distribution.

\subsection*{Metropolis--Hastings Algorithm}

\subsubsection*{Algorithm}

Given current state $x^{(t)}$:

\begin{itemize}
\item Generate proposal $Y_t \sim q(y \mid x^{(t)})$ (\emph{proposal distribution}).
\item Set
\[
X^{(t+1)} =
\begin{cases}
Y_t, & \text{with probability } \rho(x^{(t)}, Y_t), \\
x^{(t)}, & \text{with probability } 1 - \rho(x^{(t)}, Y_t),
\end{cases}
\]
where
\[
\rho(x,y)
=
\min\left\{
\frac{\pi(y)\, q(x \mid y)}{\pi(x)\, q(y \mid x)},\; 1
\right\}.
\]
\end{itemize}

\subsection*{Remarks on Metropolis--Hastings}

\begin{itemize}
\item The quantity $\rho(x,y)$ is the \emph{acceptance ratio}.
\item The resulting chain is Markov since $X^{(t+1)}$ depends only on $X^{(t)}$.
\item If $q(x \mid y) = q(y \mid x)$, the proposal is symmetric and the method reduces to the \emph{Metropolis algorithm}.
\item In this case, the acceptance probability simplifies to
\[
\frac{\pi(y_t)}{\pi(x^{(t)})}.
\]
\end{itemize}

\section{Additional Remarks}

\subsection*{Acceptance Mechanism}

If
\[
\frac{\pi(y_t)}{q(y_t \mid x^{(t)})}
>
\frac{\pi(x_t)}{q(x^{(t)} \mid y_t)},
\]
the proposal is always accepted. Otherwise, it is accepted with probability $\rho$.

\subsection*{Invariance}

The algorithm depends only on ratios:
\[
\frac{\pi(y_t)}{\pi(x^{(t)})},
\qquad
\frac{q(x^{(t)} \mid y_t)}{q(y_t \mid x^{(t)})}.
\]

Thus, normalizing constants cancel out, provided they do not depend on the conditioning variables.

\subsection*{Choice of Proposal}

The choice of proposal distribution $q$ strongly affects performance (efficiency and convergence).

A good proposal should:
\begin{itemize}
\item be easy to sample from,
\item allow easy computation of the acceptance ratio,
\item move a reasonable distance in parameter space,
\item not be rejected too frequently.
\end{itemize}

\section{Gibbs Sampling}

\subsection{Overview}

The Gibbs sampler is one of the most widely used computational methods for Bayesian inference. It was first used in statistical physics and later popularized in statistics.

The method is particularly useful for high-dimensional target distributions.

\subsection{Key Idea}

The Gibbs sampler works by \emph{sequentially sampling from univariate conditional distributions} rather than attempting to sample directly from a multivariate distribution.

\subsection{Two-Stage Gibbs Sampler}

Consider a bivariate posterior distribution
\[
\pi(\theta, \lambda \mid x).
\]

\subsubsection*{Algorithm}

\begin{itemize}
\item \textbf{Initialization:} Start with an arbitrary value $\lambda^{(0)}$.

\item \textbf{Iteration $t$:}
\begin{itemize}
\item Sample
\[
\theta^{(t)} \sim \pi_1(\theta \mid x, \lambda^{(t-1)}),
\]
\item Sample
\[
\lambda^{(t)} \sim \pi_2(\lambda \mid x, \theta^{(t)}).
\]
\end{itemize}
\end{itemize}

The joint distribution $\pi(\theta, \lambda \mid x)$ is a stationary distribution for this Markov chain.

\subsection{Multi-Stage Gibbs Sampler}

\subsubsection*{General Form}

For a joint distribution $\pi(\theta \mid x)$ with full conditionals $\pi_1, \ldots, \pi_p$, define the iterative scheme:

Given $(\theta_1^{(t)}, \ldots, \theta_p^{(t)})$,
\[
\begin{aligned}
\theta_1^{(t+1)} &\sim \pi_1(\theta_1 \mid x, \theta_2^{(t)}, \ldots, \theta_p^{(t)}), \\
\theta_2^{(t+1)} &\sim \pi_2(\theta_2 \mid x, \theta_1^{(t+1)}, \theta_3^{(t)}, \ldots, \theta_p^{(t)}), \\
&\;\;\vdots \\
\theta_p^{(t+1)} &\sim \pi_p(\theta_p \mid x, \theta_1^{(t+1)}, \ldots, \theta_{p-1}^{(t+1)}).
\end{aligned}
\]

Then
\[
\theta^{(t)} \;\to\; \theta \sim \pi(\theta \mid x).
\]

\subsection{Remarks on Gibbs Sampling}

\begin{itemize}
\item The order of updates is not essential and can be changed without affecting validity.
\item The update order can also be randomized (\emph{random scan Gibbs}).
\item It is not necessary to update every component at each iteration, provided each component is updated sufficiently often.
\item The overall process defines a Markov chain, but the individual component sequences $\theta_i^{(t)}$ are generally not Markov.
\end{itemize}

\subsection{Convergence Diagnostics}

An MCMC algorithm is said to have converged at iteration $T$ if for all $t > T$, the samples can be regarded as drawn from the stationary distribution.

\begin{itemize}
\item Lack of convergence leads to biased and unreliable estimates.
\item Diagnostics are used to detect signs of non-convergence.
\item Passing a diagnostic does \emph{not} guarantee stationarity.
\end{itemize}

\subsubsection*{Common Diagnostics}

\begin{itemize}
\item Geweke
\item Gelman and Rubin
\item Heidelberger and Welch
\item Raftery and Lewis
\end{itemize}

See the \emph{R package CODA} for implementation.

\subsubsection*{Quick Checks of Convergence}

\begin{itemize}
\item A \emph{trace plot} of $\theta^{(t)}$ against $t$ provides a basic visual diagnostic.
\item Chains may appear to converge if trapped near a local mode.
\item Running multiple chains from different starting values helps assess convergence.
\end{itemize}

\section{Autocorrelation}

\subsection*{Mixing}

The speed at which the chain explores the parameter space is called \emph{mixing}.

High autocorrelation implies slow mixing and slow convergence.

\subsection*{Autocorrelation Function}

For a sequence $\{\theta^{(1)}, \ldots, \theta^{(T)}\}$, the lag-$t$ autocorrelation function is
\[
\mathrm{acf}_t(\theta)
=
\frac{
\frac{1}{T-t} \sum_{j=1}^{T-t} (\theta^{(j)} - \bar{\theta})(\theta^{(j+t)} - \bar{\theta})
}{
\frac{1}{T-t} \sum_{j=1}^{T} (\theta^{(j)} - \bar{\theta})^2
}.
\]

\subsection*{Remarks}

\begin{itemize}
\item Gibbs sampling offers limited control over correlation structure.
\item In Metropolis-type algorithms, correlation can be adjusted through the proposal distribution.
\item Adaptive variants modify the proposal to improve mixing.
\end{itemize}

\subsection*{Practical Implications}

\begin{itemize}
\item High autocorrelation indicates slow convergence.
\item Thinning may be used to reduce correlation between retained samples.
\end{itemize}

\section{Reading: Mixture Models and Bayesian Computation}

\subsection*{Motivation for Mixtures}

Mixture models provide a flexible framework for modelling complex data distributions:
\[
f(x) = \sum_{i=1}^k p_i f(x \mid \theta_i), \quad \sum_{i=1}^k p_i = 1.
\]

They serve two key roles:

\begin{itemize}
\item \emph{Latent structure modelling}: data arise from multiple subpopulations (clustering interpretation)
\item \emph{Approximation}: mixtures act as flexible approximations to unknown densities
\end{itemize}

Thus, mixtures bridge parametric and nonparametric modeling.  

\subsection*{Latent Variable Representation}

Introduce allocation variables
\[
Z_i \sim \text{Multinomial}(1; p_1,\dots,p_k), \quad X_i \mid Z_i = j \sim f(x \mid \theta_j).
\]

This transforms the model into a \emph{complete-data formulation}, which is critical for computation.

\subsection*{Computational Challenge}

The likelihood for $n$ observations is
\[
L(\theta,p \mid x) = \prod_{i=1}^n \sum_{j=1}^k p_j f(x_i \mid \theta_j).
\]

Expanding this produces $k^n$ terms, making direct inference intractable.  

\textbf{Implication:} Closed-form Bayesian inference is generally impossible.

\subsection*{Combinatorial Explosion}

Posterior inference can be written as
\[
\pi(\theta,p \mid x)
=
\sum_{z} \omega(z)\,\pi(\theta,p \mid x,z),
\]

where $z$ are allocation vectors.

\begin{itemize}
\item Number of allocations: $k^n$
\item Only a \emph{tiny subset} have non-negligible weight
\end{itemize}

This motivates stochastic approximation methods (MCMC).  

\subsection*{EM vs Bayesian Computation}

The EM algorithm alternates:

\begin{itemize}
\item E-step: compute expected allocations
\item M-step: maximize expected complete likelihood
\end{itemize}

While efficient for MLE, EM:
\begin{itemize}
\item converges only to \emph{local modes}
\item is sensitive to initialization
\end{itemize}

\textbf{Key contrast:}
\begin{itemize}
\item EM: optimization
\item MCMC: full posterior exploration
\end{itemize}

\subsection*{Mixtures as Ill-Posed Problems}

Mixture inference is inherently unstable:

\begin{itemize}
\item Some components may receive \emph{no data}
\item Likelihood can become \emph{unbounded}
\item Small data changes $\Rightarrow$ large parameter changes
\end{itemize}

This is a classic \emph{inverse problem}.  

\subsection*{Label Switching and Identifiability}

The likelihood is invariant under permutation:
\[
(\theta_1,\dots,\theta_k) \sim (\theta_{\sigma(1)},\dots,\theta_{\sigma(k)}).
\]

Consequences:

\begin{itemize}
\item Posterior has $\mathcal{O}(k!)$ symmetric modes
\item Marginal posteriors for components are identical
\item Standard estimators (means) are meaningless
\end{itemize}

This is known as the \emph{label switching problem}.  

\subsection*{Prior Design in Mixtures}

Improper priors are problematic:

\[
\pi(\theta) = \prod_i \pi_i(\theta_i)
\quad \Rightarrow \quad
\text{posterior may be improper}.
\]

Reason:
\begin{itemize}
\item Some components may have no observations
\item Their posterior remains equal to the prior
\end{itemize}

\textbf{Solution:}
\begin{itemize}
\item Use proper priors
\item Introduce dependence between components
\end{itemize}


\subsection*{Mixtures and Nonparametric Bayes}

Mixtures naturally connect to nonparametric Bayesian methods:

\begin{itemize}
\item Kernel density estimation:
\[
\hat f(x) = \frac{1}{n} \sum_{i=1}^n K(x - x_i)
\]
is a mixture model

\item Dirichlet process:
\[
F \sim \text{DP}(F_0, \alpha)
\]
induces an \emph{infinite mixture representation}

\item Bernstein polynomials:
\[
f(x) = \sum_k p_k \text{Beta}(\alpha_k,\beta_k)
\]
approximate arbitrary densities
\end{itemize}

Thus, Bayesian nonparametrics can be viewed as \emph{mixtures with infinitely many components}.  

\subsection*{Key Takeaways for MCMC}

\begin{itemize}
\item Mixture models are computationally intractable in closed form
\item Latent variables enable tractable conditional structure
\item MCMC (Gibbs, MH) exploits this structure
\item Posterior geometry is multimodal and symmetric
\item Careful diagnostics and interpretation are required
\end{itemize}

\chapter{Dirichlet Process}

\section*{Bayesian Models Revisited}

A Bayesian statistical model consists of:

\begin{itemize}
\item a likelihood $f(y \mid \theta)$,
\item a prior distribution $\pi(\theta)$.
\end{itemize}

Uncertainty about parameters $\theta$ is modeled through the prior distribution.

Inference is based on the posterior distribution
\[
\pi(\theta \mid y)
=
\frac{f(y \mid \theta)\pi(\theta)}{\int f(y \mid \theta)\pi(\theta)\, d\theta}.
\]

\section{Nonparametric Bayesian Priors}

\subsection*{Modeling Idea}

A \emph{nonparametric Bayesian model} places a distribution on distributions.

\[
\theta \mid G \sim G
\]

where $G$ itself is random:
\[
G \sim P(G).
\]

Thus, uncertainty is modeled at two levels:
\begin{itemize}
\item uncertainty in $\theta$,
\item uncertainty in the distribution $G$.
\end{itemize}

Inference may target both $\theta$ and $G$.

\section{Finite Mixture Model}

\subsection*{Generative Model (Two Components)}

Consider a Gaussian mixture model with $K=2$:

\begin{itemize}
\item Latent assignments:
\[
z_n \stackrel{iid}{\sim} \mathrm{Categorical}(\rho_1,\rho_2),
\]

\item Observations:
\[
x_n \mid z_n \sim \mathcal{N}(\mu_{z_n}, \Sigma),
\]

\item Component means:
\[
\mu_k \stackrel{iid}{\sim} \mathcal{N}(\mu_0, \Sigma_0),
\]

\item Mixing weights:
\[
\rho_1 \sim \mathrm{Beta}(a_1,a_2), \quad \rho_2 = 1 - \rho_1.
\]
\end{itemize}

\subsection*{Inference Goal}

\begin{itemize}
\item assign data points to clusters,
\item estimate cluster parameters.
\end{itemize}

\section{Beta Distribution Review}

The Beta distribution is
\[
\mathrm{Beta}(\rho_1 \mid a_1,a_2)
=
\frac{\Gamma(a_1 + a_2)}{\Gamma(a_1)\Gamma(a_2)}
\rho_1^{a_1 - 1}(1 - \rho_1)^{a_2 - 1}.
\]

\subsection*{Gamma Function}

\begin{itemize}
\item For integer $m$: $\Gamma(m+1) = m!$
\item For $x > 0$: $\Gamma(x+1) = x\Gamma(x)$
\end{itemize}

\subsection*{Behavior}

\begin{itemize}
\item $a_1 = a_2 \to 0$: mass near boundaries
\item $a_1 = a_2 \to \infty$: concentration around $1/2$
\item $a_1 > a_2$: skew toward 1
\end{itemize}

\subsection*{Conjugacy}

If
\[
\rho_1 \sim \mathrm{Beta}(a_1,a_2), \quad z \sim \mathrm{Categorical}(\rho_1,\rho_2),
\]
then
\[
\rho_1 \mid z \sim \mathrm{Beta}(a_1 + \mathbf{1}\{z=1\},\; a_2 + \mathbf{1}\{z=2\}).
\]

\section{Dirichlet Distribution}

\subsection*{Definition}

For $\rho_{1:K}$,
\[
\mathrm{Dirichlet}(\rho_{1:K} \mid a_{1:K})
=
\frac{\Gamma\left(\sum_{k=1}^K a_k\right)}{\prod_{k=1}^K \Gamma(a_k)}
\prod_{k=1}^K \rho_k^{a_k - 1},
\]

with
\[
\rho_k > 0, \quad \sum_{k=1}^K \rho_k = 1.
\]

\subsection*{Construction via Gamma}

Let
\[
Z_j \sim \mathrm{Gamma}(a_j,1), \quad j=1,\dots,K,
\]
and define
\[
Y_j = \frac{Z_j}{\sum_{l=1}^K Z_l}.
\]

Then
\[
(Y_1,\dots,Y_K) \sim \mathrm{Dirichlet}(a_1,\dots,a_K).
\]

\subsection*{Special Case}

For $K=2$:
\[
\mathrm{Dirichlet}(a_1,a_2) = \mathrm{Beta}(a_1,a_2).
\]

\subsection*{Moments}

\[
E(Y_j) = \frac{a_j}{\sum_{l=1}^K a_l},
\quad
E(Y_j^2) = \frac{a_j(a_j+1)}{(\sum_l a_l)(1+\sum_l a_l)}.
\]

For $i \neq j$:
\[
E(Y_i Y_j) = \frac{a_i a_j}{(\sum_l a_l)(1+\sum_l a_l)}.
\]

\subsection*{Marginals}

\[
Y_j \sim \mathrm{Beta}\left(a_j,\; \sum_{i \neq j} a_i\right).
\]

\subsection*{Conjugacy}

If
\[
\rho_{1:K} \sim \mathrm{Dirichlet}(a_{1:K}), \quad z \sim \mathrm{Categorical}(\rho_{1:K}),
\]
then
\[
\rho_{1:K} \mid z \sim \mathrm{Dirichlet}(a_1',\dots,a_K'),
\quad
a_k' = a_k + \mathbf{1}\{z = k\}.
\]

\section{Finite Mixture Model (General $K$)}

\begin{itemize}
\item Mixing weights:
\[
\rho_{1:K} \sim \mathrm{Dirichlet}(a_{1:K}),
\]

\item Component means:
\[
\mu_k \sim \mathcal{N}(\mu_0,\Sigma_0),
\]

\item Assignments:
\[
z_n \sim \mathrm{Categorical}(\rho_{1:K}),
\]

\item Observations:
\[
x_n \mid z_n \sim \mathcal{N}(\mu_{z_n}, \Sigma).
\]
\end{itemize}

\section{When $K > N$}

\subsection*{Interpretation}

\begin{itemize}
\item $K$: number of possible latent groups (components)
\item clusters: number of components actually represented in the data
\end{itemize}

\subsection*{Key Observations}

\begin{itemize}
\item Number of clusters is random and typically less than $K$
\item Number of clusters grows with sample size $N$
\end{itemize}

\subsection*{Motivation}

Choosing a finite $K$ is difficult:
\begin{itemize}
\item $K$ is unknown in practice
\item inference depends heavily on $K$
\item problematic for streaming or growing datasets
\end{itemize}

\section{Choosing $K = \infty$}

\subsection*{Motivation}

In mixture models, choosing a finite $K$ is difficult:

\begin{itemize}
\item $K$ is unknown a priori,
\item inference depends strongly on $K$,
\item problematic for streaming or growing datasets.
\end{itemize}

This motivates considering the limit
\[
K = \infty.
\]

\subsection*{Infinite Mixing Weights}

We seek a sequence of strictly positive weights
\[
\rho = (\rho_1, \rho_2, \ldots)
\quad \text{such that} \quad
\sum_{k=1}^\infty \rho_k = 1.
\]

\subsection*{Dirichlet Decomposition}

From the finite Dirichlet:
\[
\rho_{1:K} \sim \mathrm{Dirichlet}(a_{1:K}),
\]

we have the decomposition:
\[
\rho_1 \sim \mathrm{Beta}\Big(a_1,\; \sum_{k=2}^K a_k\Big),
\]

and conditionally,
\[
\frac{(\rho_2,\ldots,\rho_K)}{1 - \rho_1}
\sim \mathrm{Dirichlet}(a_2,\ldots,a_K).
\]

This suggests a recursive construction.

\section{Stick-Breaking Construction}

\subsection{Finite Construction}

Define:
\[
V_1 \sim \mathrm{Beta}(a_1, a_2 + a_3 + a_4), \quad \rho_1 = V_1,
\]
\[
V_2 \sim \mathrm{Beta}(a_2, a_3 + a_4), \quad \rho_2 = (1 - V_1)V_2,
\]
\[
V_3 \sim \mathrm{Beta}(a_3, a_4), \quad \rho_3 = (1 - V_1)(1 - V_2)V_3,
\]

and so on, with
\[
\rho_K = 1 - \sum_{k=1}^{K-1} \rho_k.
\]

This is the \emph{stick-breaking} representation.

\subsection{Infinite Stick-Breaking}

Taking $K \to \infty$, define
\[
V_k \sim \mathrm{Beta}(a_k, b_k),
\]

and construct weights:
\[
\rho_k = \left(\prod_{j=1}^{k-1} (1 - V_j)\right) V_k.
\]

\section{GEM Distribution}

A special case of stick-breaking is:

\[
a_k = 1, \quad b_k = \alpha > 0.
\]

Then
\[
\rho = (\rho_1, \rho_2, \ldots) \sim \mathrm{GEM}(\alpha),
\]

with
\[
V_k \sim \mathrm{Beta}(1, \alpha),
\quad
\rho_k = \left(\prod_{j=1}^{k-1} (1 - V_j)\right) V_k.
\]

This defines a random probability distribution over a countably infinite set.

\section{Interpretation}

\subsection*{Hierarchy of Distributions}

\begin{itemize}
\item Beta: random distribution over $\{1,2\}$
\item Dirichlet: random distribution over $\{1,\ldots,K\}$
\item GEM / stick-breaking: random distribution over $\{1,2,\ldots\}$
\end{itemize}

\subsection*{Key Properties}

\begin{itemize}
\item Infinitely many components
\item Only finitely many used for finite data
\item Number of active components grows with $N$
\end{itemize}

\subsection*{Conceptual Insight}

\begin{itemize}
\item Parameters: infinite mixture components
\item Effective parameters: number of occupied clusters
\end{itemize}

\section{Readings: Foundations of the Dirichlet Process}

\subsection*{Dirichlet Process as a Distribution over Distributions}

A Dirichlet process defines a \emph{random probability measure} $P$ on a measurable space $(\mathcal{X}, \mathcal{B})$.

\begin{itemize}
\item $P$ itself is the parameter of interest
\item $P$ takes values in the space of all probability measures
\end{itemize}

The key defining property is:

For any measurable partition $(B_1, \dots, B_k)$,
\[
\big(P(B_1), \dots, P(B_k)\big)
\sim \mathrm{Dirichlet}\big(a(B_1), \dots, a(B_k)\big),
\]

where $a$ is a finite base measure.

\subsection*{Interpretation of the Base Measure}

Let
\[
a(\mathcal{X}) = \alpha, \quad \beta(\cdot) = \frac{a(\cdot)}{\alpha}.
\]

Then:

\begin{itemize}
\item $\beta$ is the \emph{base distribution} (mean measure)
\item $\alpha$ controls concentration
\end{itemize}

Key intuition:
\[
E[P(B)] = \beta(B)
\]

Thus, the Dirichlet process is centered around $\beta$.

\subsection*{Discrete Nature of the Dirichlet Process}

A fundamental property:

\begin{itemize}
\item With probability 1, $P$ is \emph{discrete}
\end{itemize}

That is,
\[
P = \sum_{n=1}^{\infty} p_n \delta_{Y_n}
\]

even if $\beta$ is continuous.

This is crucial:
\begin{itemize}
\item induces clustering
\item explains repeated parameter values
\end{itemize}

\subsection*{Posterior Conjugacy}

If
\[
P \sim \mathrm{DP}(a), \quad X \mid P \sim P,
\]

then the posterior is:
\[
P \mid X \sim \mathrm{DP}(a + \delta_X).
\]

For data $(X_1, \dots, X_n)$:
\[
P \mid X_{1:n} \sim \mathrm{DP}\left(a + \sum_{i=1}^n \delta_{X_i}\right).
\]

This is the infinite-dimensional analogue of Dirichlet conjugacy.

\subsection*{Constructive Representation (Stick-Breaking)}

A constructive definition of $P$ is:

\[
P = \sum_{n=1}^{\infty} p_n \delta_{Y_n}
\]

where:
\begin{itemize}
\item $Y_n \sim \beta$
\item weights are given by stick-breaking:
\[
p_1 = V_1,\quad
p_n = V_n \prod_{j=1}^{n-1} (1 - V_j)
\]
\item
\[
V_n \sim \mathrm{Beta}(1, \alpha)
\]
\end{itemize}

This gives an explicit generative construction of the random measure.

\subsection*{Distributional Fixed-Point Property}

The Dirichlet process satisfies the distributional equation:
\[
P = \theta \delta_Y + (1 - \theta) P',
\]

where:
\begin{itemize}
\item $\theta \sim \mathrm{Beta}(1, \alpha)$
\item $Y \sim \beta$
\item $P' \stackrel{d}{=} P$ and independent
\end{itemize}

This recursive structure underlies stick-breaking.

\subsection*{Connection to Bayesian Nonparametrics}

The Dirichlet process enables:

\begin{itemize}
\item modeling unknown distributions directly
\item infinite mixture models
\item automatic complexity control
\end{itemize}

Instead of fixing $K$, we obtain:
\[
K = \infty \quad \text{but effective clusters } < \infty.
\]

\subsection*{Practical Insight}

\begin{itemize}
\item Only finitely many atoms receive non-negligible mass
\item Sampling requires only a finite truncation
\item Posterior updates remain tractable
\end{itemize}

This resolves the key issue:
\[
\text{Unknown number of clusters} \Rightarrow \text{handled probabilistically}.
\]

\subsection*{Distributions Hierarchy Interpretation}

We can view the progression:

\begin{itemize}
\item $\mathrm{Beta}$: random distribution over $\{1,2\}$
\item $\mathrm{Dirichlet}$: random distribution over $\{1,2,\dots,K\}$
\item GEM / stick-breaking: random distribution over $\{1,2,\dots\}$
\item Dirichlet process: random distribution over $\Phi$
\end{itemize}

Formally, the Dirichlet process admits the representation:
\[
\rho = (\rho_1, \rho_2, \dots) \sim \mathrm{GEM}(\alpha),
\quad
\phi_k \overset{i.i.d.}{\sim} G_0,
\]
\[
G = \sum_{k=1}^{\infty} \rho_k \delta_{\phi_k}.
\]

Thus:
\begin{itemize}
\item $\rho_k$ controls cluster weights
\item $\phi_k$ are atom locations
\item $G$ is a discrete random measure
\end{itemize}

\subsection{Dirichlet Process Mixture Model (DPMM)}

\subsubsection*{Gaussian mixture case}

\begin{itemize}
\item Weights:
\[
\rho \sim \mathrm{GEM}(\alpha)
\]

\item Cluster parameters:
\[
\mu_k \overset{i.i.d.}{\sim} \mathcal{N}(\mu_0, \Sigma_0)
\]

\item Random measure:
\[
G = \sum_{k=1}^{\infty} \rho_k \delta_{\mu_k} \equiv DP(\alpha, \mathcal{N}(\mu_0,\Sigma_0))
\]
\end{itemize}

Latent structure:
\[
z_n \overset{i.i.d.}{\sim} \mathrm{Categorical}(\rho),
\quad
\mu_n^* = \mu_{z_n},
\quad
\mu_n^* \overset{i.i.d.}{\sim} G
\]

Observations:
\[
x_n \mid \mu_n^* \sim \mathcal{N}(\mu_n^*, \Sigma)
\]

\subsubsection*{General formulation}

More generally:

\[
\rho \sim \mathrm{GEM}(\alpha),
\quad
\phi_k \overset{i.i.d.}{\sim} G_0
\]

\[
G = \sum_{k=1}^{\infty} \rho_k \delta_{\phi_k} \equiv DP(\alpha, G_0)
\]

\[
z_n \sim \mathrm{Categorical}(\rho),
\quad
\theta_n = \phi_{z_n},
\quad
\theta_n \overset{i.i.d.}{\sim} G
\]

\[
x_n \mid \theta_n \sim F(\theta_n)
\]

\subsection*{Dirichlet Process Definition}

Let $(\mathcal{Y}, \mathcal{B}(\mathcal{Y}))$ be a measurable space.

\begin{theorem}(Dirichlet Process)
A random probability measure $F$ satisfies
\[
F \sim DP(\alpha, H)
\]
if for any measurable partition $(B_1,\dots,B_n)$,
\[
(F(B_1),\dots,F(B_n)) \sim \mathrm{Dirichlet}(\alpha H(B_1),\dots,\alpha H(B_n)).
\]
\end{theorem}

Key implication:
\[
p(F) \text{ is well-defined via Kolmogorov consistency.}
\]

\subsection*{Properties of the Dirichlet Process}

\subsubsection*{Moments}

For any measurable $A$:
\[
E[F(A)] = H(A),
\quad
\mathrm{Var}(F(A)) = \frac{H(A)(1-H(A))}{1+\alpha}.
\]

\subsubsection*{Self-similarity}

For $B$ such that $H(B)>0$:
\[
F_B \sim DP(\alpha, H_B),
\]
where $F_B$ is the restriction of $F$ to $B$.

\subsubsection*{Conjugacy}

If
\[
y_i \overset{i.i.d.}{\sim} F, \quad F \sim DP(\alpha, H),
\]
then
\[
F \mid y \sim DP\left(\alpha + n,\; \frac{1}{\alpha+n}\sum_{i=1}^n \delta_{y_i} + \frac{\alpha}{\alpha+n} H \right).
\]

In particular:
\[
E[F(A)\mid y] =
\frac{1}{\alpha+n}\sum_{i=1}^n \delta_{y_i}(A)
+
\frac{\alpha}{\alpha+n} H(A).
\]

\subsubsection*{Richness}

\begin{itemize}
\item Full support over distributions absolutely continuous w.r.t. $H$
\item Can approximate arbitrary distributions
\end{itemize}

\subsection*{Stick-Breaking Theorem}

\begin{theorem}(Sethuraman)
Let
\[
\theta_k \overset{i.i.d.}{\sim} H, \quad V_k \overset{i.i.d.}{\sim} \mathrm{Beta}(1,\alpha),
\]
and define
\[
\pi_k = V_k \prod_{l=1}^{k-1} (1 - V_l).
\]

Then
\[
\sum_{k=1}^{\infty} \pi_k \delta_{\theta_k} \sim DP(\alpha, H).
\]
\end{theorem}

\subsection*{Posterior Stick-Breaking Interpretation}

Given
\[
G \sim DP(\alpha,H), \quad \theta \sim G,
\]

we obtain:
\[
G \mid \theta \sim DP\left(\alpha+1,\; \frac{\alpha H + \delta_\theta}{\alpha+1}\right).
\]

This implies a decomposition:
\[
G = V \delta_\theta + (1-V) G',
\quad
V \sim \mathrm{Beta}(1,\alpha),
\]

where $G'$ is the residual measure.

\subsection*{Recursive Property of the DP}

For a partition $(\theta, A_1,\dots,A_K)$:

\[
(G(\theta), G(A_1), \dots, G(A_K))
\sim
\mathrm{Dirichlet}(1, \alpha H(A_1), \dots, \alpha H(A_K)).
\]

This implies:
\[
G' \sim DP(\alpha, H),
\]

showing the \emph{self-similar recursive structure} of the Dirichlet process.

\subsection*{Key Insight}

The Dirichlet process can be understood as:

\begin{itemize}
\item a distribution over discrete measures
\item constructed via stick-breaking
\item inducing clustering through shared atoms
\item recursively decomposable via Beta splitting
\end{itemize}

This structure is what enables:

\[
\text{infinite mixture models} \quad \Longleftrightarrow \quad \text{finite data-driven clustering}.
\]

\subsection*{Stick-Breaking Expansion}

We can recursively expand the Dirichlet process draw:
\[
G \sim DP(\alpha, H)
\]

\[
G = V_1 \delta_{\theta_1^*} + (1 - V_1) G_1
\]

\[
G = V_1 \delta_{\theta_1^*} + (1 - V_1)\big(V_2 \delta_{\theta_2^*} + (1 - V_2) G_2\big)
\]

Continuing:
\[
G = \sum_{k=1}^{\infty} \pi_k \delta_{\theta_k^*},
\]

where
\[
\pi_k = V_k \prod_{l=1}^{k-1} (1 - V_l),
\quad
V_k \sim \mathrm{Beta}(1,\alpha),
\quad
\theta_k^* \sim H.
\]

\subsection*{Interpretation}

\begin{itemize}
\item Draws from a DP are \emph{discrete} measures
\item Mass is allocated via stick-breaking
\item Atoms $\theta_k^*$ are drawn from the base distribution $H$
\end{itemize}

\subsection*{Marginalizing the Weights}

Consider a simple case:
\[
\rho_1 \sim \mathrm{Beta}(a_1,a_2), \quad z_n \sim \mathrm{Categorical}(\rho_1,\rho_2).
\]

Integrating out $\rho_1$, we obtain:
\[
p(z_n = 1 \mid z_{1:n-1})
=
\frac{a_{1,n}}{a_{1,n} + a_{2,n}},
\]

where
\[
a_{1,n} = a_1 + \sum_{m=1}^{n-1} \mathbf{1}\{z_m = 1\},
\quad
a_{2,n} = a_2 + \sum_{m=1}^{n-1} \mathbf{1}\{z_m = 2\}.
\]

\subsection{Pólya Urn Interpretation}

This predictive rule corresponds to a Pólya urn:

\begin{itemize}
\item Choose a ball with probability proportional to its count
\item Replace it and add another ball of the same color
\end{itemize}

Then:
\[
\lim_{n \to \infty}
\frac{\# \text{ orange}}{\# \text{ total}}
\stackrel{d}{=} \mathrm{Beta}(a_{\text{orange}}, a_{\text{green}}).
\]

\subsection*{Multivariate Extension}

For
\[
\rho_{1:K} \sim \mathrm{Dirichlet}(a_{1:K}), \quad z_n \sim \mathrm{Categorical}(\rho_{1:K}),
\]

we obtain:
\[
p(z_n = k \mid z_{1:n-1})
=
\frac{a_{k,n}}{\sum_{j=1}^K a_{j,n}},
\]

where
\[
a_{k,n} = a_k + \sum_{m=1}^{n-1} \mathbf{1}\{z_m = k\}.
\]

This corresponds to a \emph{multivariate Pólya urn}.

\subsection{Blackwell--MacQueen Pólya Urn}

We now consider the infinite limit.

Let:
\[
G \sim DP(\alpha, H), \quad \theta_i \mid G \sim G.
\]

Then:

\subsubsection*{First sample}

\[
\theta_1 \sim H,
\quad
G \mid \theta_1 \sim DP\left(\alpha + 1, \frac{\alpha H + \delta_{\theta_1}}{\alpha + 1}\right).
\]

\subsubsection*{Second sample}

\[
\theta_2 \mid \theta_1
\sim
\frac{\alpha H + \delta_{\theta_1}}{\alpha + 1}.
\]

\subsubsection*{General step}

\[
\theta_n \mid \theta_{1:n-1}
\sim
\frac{\alpha H + \sum_{i=1}^{n-1} \delta_{\theta_i}}{\alpha + n - 1}.
\]

\subsection*{Urn Interpretation}

\begin{itemize}
\item With probability proportional to $\alpha$, draw a new value from $H$
\item With probability proportional to existing counts, reuse a previous value
\end{itemize}

Equivalently:

\begin{itemize}
\item Draw a ball proportional to its mass
\item If it is \emph{new}, create a new color
\item Otherwise, reinforce the existing color
\end{itemize}

\subsection*{Key Insight}

This produces a sequence:
\[
\theta_1, \theta_2, \dots
\]

with:
\begin{itemize}
\item clustering (shared values)
\item increasing number of clusters
\item probability of new cluster controlled by $\alpha$
\end{itemize}

\section{Chinese Restaurant Process (CRP)}

\subsection*{Construction}

The Chinese Restaurant Process provides an equivalent formulation of the Dirichlet process clustering mechanism.

\begin{itemize}
\item Each observation (customer) arrives sequentially.
\item Customer $n$:
\begin{itemize}
\item sits at an existing table $k$ with probability proportional to the number of customers already there,
\item starts a new table with probability proportional to $\alpha$.
\end{itemize}
\end{itemize}

Formally:
\[
p(z_n = k \mid z_{1:n-1}) =
\begin{cases}
\frac{n_k}{n-1+\alpha}, & \text{existing table } k, \\
\frac{\alpha}{n-1+\alpha}, & \text{new table}.
\end{cases}
\]

\subsection*{Partition Interpretation}

Let $\Pi_N$ denote the partition of $\{1,\dots,N\}$ induced by assignments $z_1,\dots,z_N$.

Example:
\[
z_1 = z_2 = z_7 = z_8 = 1,\quad
z_3 = z_5 = z_6 = 2,\quad
z_4 = 3
\]
\[
\Rightarrow
\Pi_8 = \{\{1,2,7,8\}, \{3,5,6\}, \{4\}\}.
\]

A partition consists of mutually exclusive and exhaustive subsets of $\{1,\dots,N\}$.

\subsection*{Probability of a Seating Arrangement}

Let $K_N$ be the number of tables and let $C \in \Pi_N$ denote a cluster.

Then:
\[
\mathbb{P}(\Pi_N = \pi_N)
=
\frac{\alpha^{K_N} \prod_{C \in \Pi_N} (|C|-1)!}{\alpha(\alpha+1)\cdots(\alpha+N-1)}.
\]

\subsection*{Exchangeability}

The probability of a partition does not depend on the order of customers:
\[
\mathbb{P}(\Pi_N = \pi_N)
\text{ is invariant under permutations.}
\]

Thus, the CRP is \emph{exchangeable}.

\subsection*{Conditional Assignment Rule}

Given $\Pi_{N,-n}$ (partition without customer $n$):
\[
p(\Pi_N \mid \Pi_{N,-n}) =
\begin{cases}
\frac{|C|}{\alpha+N-1}, & \text{if } n \text{ joins cluster } C, \\
\frac{\alpha}{\alpha+N-1}, & \text{if } n \text{ creates new cluster}.
\end{cases}
\]

\section*{Representations of the Dirichlet Process}

The Dirichlet process admits several equivalent representations:

\begin{itemize}
\item \emph{Posterior form:}
\[
G \sim DP(\alpha,H), \quad \theta \mid G \sim G
\]
\[
\Longleftrightarrow
\theta \sim H,\quad
G \mid \theta \sim DP\left(\alpha+1, \frac{\alpha H + \delta_\theta}{\alpha+1}\right)
\]

\item \emph{Pólya urn:}
\[
\theta_n \mid \theta_{1:n-1}
\sim
\frac{\alpha H + \sum_{i=1}^{n-1} \delta_{\theta_i}}{\alpha+n-1}
\]

\item \emph{CRP:}
\[
p(z_n = k \mid z_{1:n-1}) =
\begin{cases}
\frac{n_k}{n-1+\alpha}, \\
\frac{\alpha}{n-1+\alpha}
\end{cases}
\]

\item \emph{Stick-breaking:}
\[
\pi_k = \beta_k \prod_{i=1}^{k-1}(1-\beta_i),\quad
\beta_k \sim \mathrm{Beta}(1,\alpha)
\]
\[
G = \sum_{k=1}^{\infty} \pi_k \delta_{\theta_k^*}
\]
\end{itemize}

\section*{Exchangeability and Joint Distribution}

From the predictive rule:
\[
\theta_n \mid \theta_{1:n-1}
\sim
\frac{\alpha H + \sum_{i=1}^{n-1} \delta_{\theta_i}}{\alpha+n-1},
\]

the joint distribution is:
\[
p(\theta_1,\dots,\theta_n)
=
\frac{\alpha^K \prod_{k=1}^K h(\theta_k^*) (m_k-1)!}{\prod_{i=1}^n (\alpha + i - 1)},
\]

where:
\begin{itemize}
\item $\theta_k^*$ are distinct values,
\item $m_k$ are their multiplicities,
\item $h(\theta)$ is the density under $H$.
\end{itemize}

This distribution is exchangeable.

By De Finetti’s theorem, this implies the existence of a random measure $G$ such that
\[
\theta_1,\theta_2,\dots \overset{iid}{\sim} G,
\quad G \sim DP(\alpha,H).
\]

\section{Number of Clusters}

Let $K_n$ denote the number of clusters among $n$ observations.

\subsection{Expected Value}

Define indicator:
\[
I_i = \mathbf{1}\{\text{new cluster at step } i\}.
\]

Then:
\[
\mathbb{E}[K_n]
=
\sum_{i=1}^n \mathbb{E}[I_i]
=
\sum_{i=1}^n \frac{\alpha}{\alpha+i-1}.
\]

Thus:
\[
\mathbb{E}[K_n] = O(\alpha \log n).
\]

\subsection{Variance}

\[
\mathrm{Var}(K_n)
=
\sum_{i=1}^n
\frac{\alpha(i-1)}{(\alpha+i-1)^2}.
\]

\subsection{Distribution of $K_n$}

The distribution satisfies:
\[
\mathbb{P}(K_n = k)
=
\frac{\alpha}{\alpha+n-1} \mathbb{P}(K_{n-1}=k-1)
+
\frac{n-1}{\alpha+n-1} \mathbb{P}(K_{n-1}=k).
\]

Boundary conditions:
\[
\mathbb{P}(K_1 = 1) = 1.
\]

Closed form:
\[
\mathbb{P}(K_n = k)
=
C_n(k)\,\alpha^k \frac{\Gamma(\alpha)}{\Gamma(\alpha+n)},
\]

where $C_n(k)$ are unsigned Stirling numbers of the first kind.

\subsection{Interpretation}

\begin{itemize}
\item Rich-get-richer effect: larger clusters attract more points
\item Few large clusters, many small ones
\item Number of clusters grows sublinearly in $n$
\end{itemize}

\chapter{Sampling Techniques}

\section{Sampling Techniques for DP Mixture Models}

\subsection*{Gaussian Likelihood Setup}

Assume:
\[
y \sim \mathcal{N}(\mu, \sigma^2),
\quad
p(y \mid \mu,\sigma^2) = \frac{1}{\sqrt{2\pi\sigma^2}} \exp\left(-\frac{1}{2\sigma^2}(y-\mu)^2\right).
\]

For i.i.d. data $y = (y_1,\dots,y_n)$:
\[
L(\mu,\sigma^2)
\propto
\left(\frac{1}{\sigma^2}\right)^{n/2}
\exp\left(-\frac{1}{2\sigma^2}\sum_i (y_i-\mu)^2\right).
\]

We focus on:
\[
p(\mu,\sigma^2 \mid y) = p(\mu \mid \sigma^2, y)\, p(\sigma^2 \mid y).
\]

\subsection*{Semi-Conjugate Prior}

Assume independent priors:
\[
\pi(\mu,\sigma^2) = \pi(\mu)\pi(\sigma^2),
\]

\[
\mu \sim \mathcal{N}(\mu_0,\sigma_0^2),
\quad
\sigma^2 \sim IG\left(\frac{\nu_0}{2}, \frac{\delta_0}{2}\right).
\]

Then:
\[
\mu \mid \sigma^2, y \sim \mathcal{N}(\mu_n, \tau_n^2),
\]

\[
\mu_n =
\frac{\mu_0/\sigma_0^2 + \bar{y}\, n/\sigma^2}{1/\sigma_0^2 + n/\sigma^2},
\quad
\tau_n^2 = \frac{1}{1/\sigma_0^2 + n/\sigma^2}.
\]

\[
\sigma^2 \mid y \text{ not in closed form}, \quad
\sigma^2 \mid y,\mu \sim IG.
\]

\subsection{Finite Mixture Models}

A flexible density:
\[
f(y) = \sum_{k=1}^K \tau_k f(y \mid \theta_k),
\]

where:
\[
\tau_k \in (0,1), \quad \sum_k \tau_k = 1.
\]

\subsection{Bayesian Mixture Model}

\[
y_i \mid \tau,\theta \sim \sum_{k=1}^K \tau_k f(y_i \mid \theta_k),
\]

\[
\tau \sim \mathrm{Dir}(\alpha,\dots,\alpha),
\quad
\theta_k \sim \pi_k(\theta_k),
\quad
K \sim \pi(K).
\]

\subsection{Gibbs Sampling (Finite Case)}

For fixed $K$:

\begin{itemize}
\item Assignment:
\[
P(z_i^{(t)} = k \mid \cdot)
\propto
\tau_k^{(t-1)} f(y_i \mid \theta_k^{(t-1)}).
\]

\item Weights:
\[
\tau \mid z \sim \mathrm{Dir}(\alpha_1 + \hat n_1,\dots,\alpha_K + \hat n_K),
\]

where $\hat n_k = \sum_i \mathbf{1}(z_i = k)$.

\item Parameters:
\[
p(\theta_k \mid \cdot)
\propto
\pi_k(\theta_k)\prod_{i:z_i=k} f(y_i \mid \theta_k).
\]
\end{itemize}

\section{CRP Mixture Model}

\subsection*{Generative Model}

\[
\Pi_N \sim CRP(N,\alpha),
\]

\[
\forall C \in \Pi_N,\quad \phi_C \overset{i.i.d.}{\sim} G_0,
\]

\[
x_n \mid \phi_C \sim F(\phi_C), \quad n \in C.
\]

\subsection*{Inference Goal}

\[
p(\Pi_N \mid x_{1:N}).
\]

\subsection*{Gibbs Sampler (Collapsed)}

\[
p(\Pi_N \mid \Pi_{N,-n}, x)
=
\begin{cases}
\frac{|C|}{\alpha+N-1}\, p(x_{C \cup \{n\}} \mid x_C), & \text{join } C, \\
\frac{\alpha}{\alpha+N-1}\, p(x_n), & \text{new cluster}.
\end{cases}
\]

For Gaussian case:
\[
p(x_{C \cup \{n\}} \mid x_C)
=
\mathcal{N}(\tilde m,\tilde \Sigma + \Sigma),
\]

\[
\tilde \Sigma^{-1} = \Sigma_0^{-1} + |C|\Sigma^{-1},
\quad
\tilde m = \tilde \Sigma\left(\Sigma^{-1}\sum_{m\in C} x_m + \Sigma_0^{-1}\mu_0\right).
\]

\section{DP Mixture Model (DPM)}

\[
y_i \sim \int p(y_i \mid \theta) G(d\theta),
\quad
G \sim DP(M,G_0).
\]

Example:
\[
y_i \sim \int \mathcal{N}(y_i \mid \mu,\sigma^2) G(d\mu).
\]

\section{Alternative Representations}

\subsection{Stick-Breaking Form}

\[
y_i \mid (w_h,\tilde\theta_h)
\sim
\sum_{h=1}^{\infty} w_h p(y_i \mid \tilde\theta_h),
\]

\[
\tilde\theta_h \sim G_0,
\quad
w_h = v_h \prod_{k<h}(1-v_k),
\quad
v_h \sim \mathrm{Beta}(1,M).
\]

\subsection{Clustering Representation}

\[
y_i \mid \theta_i \sim p(y_i \mid \theta_i),
\quad
\theta_i \mid G \sim G,
\quad
G \sim DP(M,G_0).
\]

Integrating out $G$:
\[
(\theta_1,\dots,\theta_n) \sim \text{CRP-induced distribution}.
\]

Predictive:
\[
p(\theta_{n+1} \mid \theta_{1:n})
\propto
\sum_{j=1}^{k_n} n_j \delta_{\theta_j^*} + M G_0.
\]

\section{Collapsed Gibbs Sampling}

\subsection{Conditional for Parameters}

\[
\theta_i \mid \theta_{-i}, y
\propto
\sum_{j=1}^{k^-} n_j^- p(y_i \mid \theta_j^*) \delta_{\theta_j^*}
+
M p(y_i \mid \theta_i) G_0(\theta_i).
\]

Posterior:
\[
p(\theta_i \mid y_i, G_0)
=
\frac{p(y_i \mid \theta_i)G_0(\theta_i)}{\int p(y_i \mid \theta)G_0(d\theta)}.
\]

\subsection{Cluster Assignment Update}

\[
p(s_i = j \mid s^{-}, y)
\propto
\begin{cases}
n_j^- \int p(y_i \mid \theta_j^*) d p(\theta_j^* \mid y_j^-), & j \le k^-, \\
M \int p(y_i \mid \theta) dG_0(\theta), & j = k^-+1.
\end{cases}
\]

\subsection{Cluster Parameter Update}

\[
p(\theta_j^* \mid s,y)
\propto
G_0(\theta_j^*) \prod_{i:s_i=j} p(y_i \mid \theta_j^*).
\]

\section{Gaussian DPM Example}

\[
p(y_i \mid \theta_i) = \mathcal{N}(\theta_i,\sigma^2),
\quad
G_0 = \mathcal{N}(m,B).
\]

\[
\int p(y_i \mid \theta)dG_0(\theta)
=
\mathcal{N}(y_i \mid m, B+\sigma^2).
\]

Cluster predictive:
\[
\int p(y_i \mid \theta_j^*) d p(\theta_j^* \mid y_j^-)
=
\mathcal{N}(y_i \mid m_j^-, V_j^- + \sigma^2),
\]

\[
V_j^- = \left(\frac{1}{B} + \frac{n_j^-}{\sigma^2}\right)^{-1},
\quad
m_j^- = V_j^- \left(\frac{m}{B} + \frac{1}{\sigma^2}\sum_{h\in S_j^-} y_h \right).
\]

Posterior:
\[
p(\theta_j^* \mid y)
=
\mathcal{N}(m_j, V_j).
\]

\section{Reading: Mixture of Dirichlet Process Models}

\subsection*{Model Representation}

The Dirichlet process mixture model is
\[
y_i \mid \theta_i \sim F(\theta_i), \quad \theta_i \sim G, \quad G \sim DP(\alpha, G_0)
\]

Integrating out $G$ gives
\[
\theta_i \mid \theta_{1:i-1}
=
\sum_{j<i} \frac{1}{i-1+\alpha}\delta_{\theta_j}
+
\frac{\alpha}{i-1+\alpha} G_0
\]

This induces clustering through repeated values of $\theta_i$.

\subsection*{Partition View}

Let $c_i$ denote cluster assignments. Then
\[
P(c_i = c \mid c_{-i})
=
\begin{cases}
\frac{n_{-i,c}}{n-1+\alpha} \\
\frac{\alpha}{n-1+\alpha}
\end{cases}
\]

and
\[
y_i \mid c_i = c \sim F(\theta_c)
\]

Thus inference reduces to learning
\[
(\Pi_n, \{\theta_c\})
\]

\subsection*{Non-Conjugate Difficulty}

Key terms required in Gibbs sampling are
\[
\int F(y_i \mid \theta)\, dG_0(\theta)
\quad \text{and} \quad
p(\theta \mid y_i) \propto F(y_i \mid \theta)G_0(\theta)
\]

These are not available in closed form in general.

\subsection*{No-Gaps Representation}

Clusters are indexed as $1,\dots,k$ with no empty labels. This avoids redundant representations of the same partition and simplifies computation.

\subsection*{Augmented Sampling Idea}

Introduce auxiliary parameters
\[
\theta_1^*, \dots, \theta_m^* \sim G_0
\]

Then approximate the new cluster term via
\[
\alpha \int F(y_i \mid \theta)\, dG_0(\theta)
\approx
\frac{\alpha}{m} \sum_{j=1}^m F(y_i \mid \theta_j^*)
\]

\subsection*{Gibbs Updates}

Cluster assignments are updated via
\[
P(c_i = c \mid \cdot)
\propto
\begin{cases}
n_{-i,c} F(y_i \mid \theta_c) \\
\frac{\alpha}{m} F(y_i \mid \theta_c^*)
\end{cases}
\]

Cluster parameters are updated via
\[
\theta_c \mid y_c \sim p(\theta_c \mid y_c)
\]

\subsection*{Main Insight}

Non-conjugacy can be handled by expanding the state space and performing Gibbs updates in the augmented model, avoiding direct evaluation of intractable integrals.



\section{Reading: Sampling Methods for Dirichlet Process Mixtures}

\subsection*{Dirichlet Process Mixture Representation}

The model is
\[
y_i \mid \theta_i \sim F(\theta_i), \quad \theta_i \sim G, \quad G \sim DP(\alpha, G_0)
\]

which implies
\[
\theta_i \mid \theta_{1:i-1}
=
\sum_{j<i} \frac{1}{i-1+\alpha}\delta_{\theta_j}
+
\frac{\alpha}{i-1+\alpha} G_0
\]

This is equivalent to a Chinese restaurant process on cluster assignments.

\subsection*{Collapsed Gibbs Sampling}

If conjugate, integrate out $\theta_c$:
\[
P(c_i = c \mid c_{-i}, y)
\propto
n_{-i,c} \cdot p(y_i \mid y_{c,-i})
\]

For a new cluster:
\[
\propto
\alpha \int F(y_i \mid \theta)\, dG_0(\theta)
\]

\subsection*{Failure Under Non-Conjugacy}

The integral
\[
\int F(y_i \mid \theta)\, dG_0(\theta)
\]
is not tractable, so collapsed Gibbs cannot be used.

\subsection*{Metropolis-Hastings Update}

Propose
\[
\theta^* \sim G_0
\]

Accept with probability
\[
\rho = \min\left(1, \frac{F(y_i \mid \theta^*)}{F(y_i \mid \theta_i)}\right)
\]

The prior cancels, so no integral is needed.

\subsection*{Auxiliary Gibbs Sampling}

Introduce auxiliary draws
\[
\theta_1^*, \dots, \theta_m^* \sim G_0
\]

Update assignments via
\[
P(c_i = c \mid \cdot)
\propto
\begin{cases}
n_{-i,c} F(y_i \mid \theta_c) \\
\frac{\alpha}{m} F(y_i \mid \theta_c^*)
\end{cases}
\]

\subsection*{Mixing Behavior}

Standard Gibbs updates one observation at a time:
\[
c_i \to c_i'
\]

This leads to slow exploration of partitions and difficulty in splitting or merging clusters.

\subsection*{Conceptual Summary}

Posterior inference is over
\[
p(\Pi_n, \{\theta_c\} \mid y)
\]

Key challenges are
\[
\text{non-conjugacy} \quad \text{and} \quad \text{infinite mixture structure}
\]

Efficient sampling requires combining
\[
\text{Gibbs}, \quad \text{Metropolis-Hastings}, \quad \text{augmentation}
\]

\section{Sampling Techniques: Slice Sampling for DP Mixtures}

\subsection*{Hyperpriors for Base Measure}

Assume the base measure is parameterized as $\phi = (m,B)$ with priors
\[
m \sim \mathcal{N}(m_0, D), \quad B \sim IGamma(a,b)
\]

Given cluster parameters $\{\theta_j^*\}_{j=1}^k$, the posterior for $m$ is
\[
m \mid \cdots \sim \mathcal{N}(m_1, D_1),
\quad
D_1^{-1} = \frac{1}{D} + \frac{k}{B},
\quad
m_1 = D_1\left(\frac{m_0}{D} + \frac{1}{B}\sum_{j=1}^k \theta_j^*\right)
\]

The posterior for $B$ is
\[
B \mid \cdots \sim IGamma(a_1, b_1),
\quad
a_1 = a + \frac{k}{2},
\quad
b_1 = b + \frac{1}{2}\sum_{j=1}^k (\theta_j^* - m)^2
\]

\subsection*{Sampling the Concentration Parameter}

We use the marginal:
\[
p(k \mid M) \propto M^k \frac{\Gamma(M)}{\Gamma(M+n)}
\]

Introduce auxiliary variable:
\[
\eta \mid M,k \sim Beta(M+1, n)
\]

If prior $M \sim Gamma(c,d)$, then:
\[
M \mid \eta,k \sim
\pi \cdot Gamma(c+k, d - \log \eta)
+
(1-\pi) \cdot Gamma(c+k-1, d - \log \eta)
\]

where mixing weight depends on $(k,n,\eta)$.

\subsection*{No-Gaps Algorithm}

Clusters are labeled consecutively:
\[
c_i \in \{1,\dots,k\}
\]

with relabeling applied whenever clusters become empty.

For each observation $i$, remove $y_i$ and define:
\[
k^- = \text{\# clusters without } i, \quad n_j^- = \sum_{\ell \neq i} \mathbf{1}(c_\ell = j)
\]

If $y_i$ is not a singleton:
\[
P(c_i = j \mid \cdot) \propto
\begin{cases}
n_j^- p(y_i \mid \theta_j^*) & j \le k^- \\
\frac{M}{k^- + 1} p(y_i \mid \theta_{k^-+1}^*) & j = k^- + 1
\end{cases}
\]

If $y_i$ is a singleton:
\[
P(\text{keep } c_i) = \frac{k^- - 1}{k^-}, \quad
P(\text{resample}) = \frac{1}{k^-}
\]

\subsection*{Updating Cluster Parameters}

Given assignments $\{c_i\}$:
\[
p(\theta_j^* \mid y) \propto G_0(\theta_j^*) \prod_{i:c_i=j} p(y_i \mid \theta_j^*)
\]

If non-conjugate, sample via Metropolis-Hastings.

New cluster candidates:
\[
\theta_{k+1}^* \sim G_0
\]

\subsection*{Slice Sampling Idea}

Target density:
\[
\pi(x) = \frac{f(x)}{Z}, \quad Z = \int f(s)\,ds
\]

Introduce auxiliary variable $u$:
\[
u \sim Uniform(0, f(x))
\]

Define joint:
\[
p(x,u) = \frac{1}{Z} \mathbf{1}(0 < u < f(x))
\]

Marginalizing:
\[
p(x) = \int_0^{f(x)} \frac{1}{Z} du = \frac{f(x)}{Z}
\]

\subsection*{Slice Sampling Updates}

Sample:
\[
u_t \sim Uniform(0, f(x_t))
\]

Define slice:
\[
S(u) = \{x : f(x) \ge u\}
\]

Then update:
\[
x_{t+1} \sim Uniform(S(u_t))
\]

\subsection*{Slice Sampling for DP Mixtures}

Likelihood:
\[
p(y_i \mid \{w_h\}, \{\tilde{\theta}_h\}) = \sum_{h=1}^\infty w_h p(y_i \mid \tilde{\theta}_h)
\]

Introduce slice variable:
\[
u_i \sim Uniform(0, w_{z_i})
\]

Augmented likelihood:
\[
p(y_i, u_i) = \sum_{h=1}^\infty \mathbf{1}(u_i < w_h) p(y_i \mid \tilde{\theta}_h)
\]

Key property:
\[
u_i > 0 \Rightarrow \text{only finitely many } w_h > u_i
\]

Thus:
\[
\sum_{h=1}^\infty \rightarrow \sum_{h: w_h > u_i}
\]

which is finite.

\subsection*{Marginalization Check}

\[
\int p(y_i, u_i) du_i
=
\sum_{h=1}^\infty p(y_i \mid \tilde{\theta}_h) \int_0^{w_h} du_i
=
\sum_{h=1}^\infty w_h p(y_i \mid \tilde{\theta}_h)
\]

Thus the original model is recovered.

\subsection*{Key Insight}

Slice sampling converts:
\[
\text{infinite mixture} \quad \rightarrow \quad \text{finite active components per iteration}
\]

by introducing auxiliary variables $\{u_i\}$.

This enables exact inference without truncation.

\subsection*{Slice Sampler: Component Indicators}

Introduce latent indicators $r_i \in \{1,2,\dots\}$ such that
\[
r_i = h \quad \Longleftrightarrow \quad y_i \text{ is generated from } \tilde{\theta}_h
\]

The augmented likelihood becomes
\[
p(y_i, u_i, r_i \mid \{w_h\}, \{\tilde{\theta}_h\})
=
\mathbf{1}(u_i < w_{r_i}) \, p(y_i \mid \tilde{\theta}_{r_i})
\]

Marginalizing out $r_i$:
\[
\sum_{h=1}^\infty \mathbf{1}(u_i < w_h) p(y_i \mid \tilde{\theta}_h)
\]
recovers the slice-augmented likelihood.

The full joint:
\[
p(y, u, r \mid \{w_h\}, \{\tilde{\theta}_h\})
=
\prod_{i=1}^n \mathbf{1}(u_i < w_{r_i}) p(y_i \mid \tilde{\theta}_{r_i})
\]

\subsection*{Slice Sampler: Conditional Updates}

\textbf{Slice variables:}
\[
u_i \mid r_i, \{w_h\} \sim Uniform(0, w_{r_i})
\]

\textbf{Component indicators:}
\[
P(r_i = h \mid u_i, y_i, \{\tilde{\theta}_h\})
\propto
\mathbf{1}(w_h > u_i)\, p(y_i \mid \tilde{\theta}_h)
\]

Only components with $w_h > u_i$ are considered.

\textbf{Component parameters:}
\[
p(\tilde{\theta}_h \mid \{y_i : r_i = h\})
\propto
G_0(\tilde{\theta}_h) \prod_{i:r_i=h} p(y_i \mid \tilde{\theta}_h)
\]

\textbf{Stick-breaking weights:}

Let
\[
w_h = v_h \prod_{k<h} (1 - v_k)
\]

Then
\[
v_h \mid \{r_i\}
\sim
Beta\left(1 + n_h,\; M + \sum_{k>h} n_k\right),
\quad
n_h = \sum_i \mathbf{1}(r_i = h)
\]

\subsection*{Finite Approximation to DP}

Instead of infinite mixtures:
\[
G = \sum_{h=1}^\infty w_h \delta_{\tilde{\theta}_h}
\]

approximate with finite truncation:
\[
G_H = \sum_{h=1}^H w_h \delta_{\tilde{\theta}_h}
\]

The marginal distributions are close:
\[
\| p_\infty - p_H \|_1 \le 4n e^{-(H-1)/M}
\]

Thus small $H$ can suffice when $M$ is moderate.

\subsection*{Finite DP Posterior Structure}

Joint posterior:
\[
p(\{r_i\}, \{\tilde{\theta}_h\}, \{v_h\} \mid y)
\propto
\prod_{i=1}^n p(y_i \mid \tilde{\theta}_{r_i})
\prod_{i=1}^n p(r_i \mid \{v_h\})
\prod_{h=1}^H dG_0(\tilde{\theta}_h)
\prod_{h=1}^{H-1} p(v_h \mid M)
\]

Conditionals:

\[
p(\tilde{\theta}_h \mid \cdot)
\propto
G_0(\tilde{\theta}_h) \prod_{i:r_i=h} p(y_i \mid \tilde{\theta}_h)
\]

\[
v_h \mid \cdot \sim Beta\left(1 + n_h,\; M + \sum_{k>h} n_k\right)
\]

\[
P(r_i = h \mid \cdot)
\propto
w_h p(y_i \mid \tilde{\theta}_h)
\]

\subsection*{Key Comparison}

\begin{itemize}
\item Slice sampler: exact, dynamically truncates via $\{u_i\}$
\item Finite DP: fixed truncation level $H$
\item Both reduce infinite mixture to finite computation per iteration
\end{itemize}

\section{Reading: Gibbs Sampling for Stick-Breaking Priors}

\subsection*{Stick-Breaking Random Measures}

A general stick-breaking prior can be written as
\[
G = \sum_{k=1}^{\infty} p_k \delta_{Z_k}
\]
where
\[
p_1 = V_1, \quad p_k = V_k \prod_{l<k}(1 - V_l), \quad V_k \sim Beta(a_k, b_k)
\]

The atoms satisfy
\[
Z_k \overset{iid}{\sim} H
\]

This defines a broad class of random probability measures including DP and Pitman–Yor.

\subsection*{Finite vs Infinite Construction}

Finite case:
\[
\sum_{k=1}^N p_k = 1 \quad \text{with } V_N = 1
\]

Infinite case requires
\[
\sum_{k=1}^\infty \mathbb{E}[\log(1 - V_k)] = -\infty
\]

to ensure weights sum to 1 almost surely.

\subsection*{Dirichlet Process Special Case}

DP corresponds to
\[
V_k \sim Beta(1, \alpha)
\]

giving
\[
G \sim DP(\alpha, H)
\]

This provides a constructive definition equivalent to Ferguson’s original formulation.

\subsection*{Generalized Gibbs Sampling}

Two main approaches:

\textbf{Pólya urn Gibbs:}
\[
\theta_{n+1} \mid \theta_{1:n}
\sim
\frac{\alpha}{\alpha+n} H + \sum_{k} \frac{n_k}{\alpha+n} \delta_{\theta_k^*}
\]

\textbf{Blocked Gibbs:}
Directly sample
\[
\{V_k\}, \{Z_k\}, \{p_k\}
\]

This avoids marginalization and improves mixing.

\subsection*{Pitman–Yor Extension}

Stick-breaking form:
\[
V_k \sim Beta(1-a, b + ka)
\]

This yields richer clustering behavior than DP.

\subsection*{Key Insight}

Stick-breaking provides:

\begin{itemize}
\item explicit construction of random measures
\item direct sampling schemes
\item flexibility beyond Dirichlet process
\end{itemize}

\section{Reading: Slice Sampling Mixture Models}

\subsection*{Dirichlet Process Mixture Model}

The model:
\[
f(y) = \int K(y \mid \phi)\, dP(\phi), \quad P \sim DP(M, P_0)
\]

Stick-breaking form:
\[
P = \sum_{j=1}^\infty w_j \delta_{\phi_j}, \quad w_j = z_j \prod_{l<j}(1-z_l), \quad z_j \sim Beta(1,M)
\]

\subsection*{Slice Sampling Idea}

Introduce auxiliary variable:
\[
u_i \sim Uniform(0, w_{d_i})
\]

Joint:
\[
p(y_i, u_i) = \sum_{j=1}^\infty \mathbf{1}(u_i < w_j) K(y_i \mid \phi_j)
\]

This restricts to finite set:
\[
A_u = \{j: w_j > u_i\}
\]

\subsection*{Latent Component Indicator}

Introduce
\[
d_i \in A_u
\]

Then
\[
p(y_i, u_i, d_i) = \mathbf{1}(u_i < w_{d_i}) K(y_i \mid \phi_{d_i})
\]

This converts infinite mixture into finite conditional problem.

\subsection*{Gibbs Updates}

\[
\phi_j \mid \cdot \propto p_0(\phi_j)\prod_{i:d_i=j} K(y_i \mid \phi_j)
\]

\[
z_j \sim Beta\left(1 + \sum_i \mathbf{1}(d_i=j),\; M + \sum_i \mathbf{1}(d_i > j)\right)
\]

\[
u_i \sim Uniform(0, w_{d_i})
\]

\[
P(d_i = k) \propto \mathbf{1}(w_k > u_i) K(y_i \mid \phi_k)
\]

\subsection*{Key Insight}

Slice variables ensure:

\begin{itemize}
\item only finitely many components active per iteration
\item exact inference without truncation
\item easy extension to general stick-breaking priors
\end{itemize}

\section{Reading: Bayesian Density Estimation via DP Mixtures}

\subsection*{Model Setup}

Data:
\[
y_i \mid \theta_i \sim N(\mu_i, V_i)
\]

Prior:
\[
\theta_i \sim G, \quad G \sim DP(\alpha G_0)
\]

\subsection*{Clustering Property}

Conditional distribution:
\[
\theta_i \mid \theta_{-i}
\sim
\frac{\alpha}{\alpha+n-1} G_0 + \sum_{j\neq i} \frac{1}{\alpha+n-1} \delta_{\theta_j}
\]

Leads to clustering (ties among $\theta_i$).

\subsection*{Posterior Predictive}

\[
y_{n+1} \mid y_{1:n}
\sim
\frac{\alpha}{\alpha+n} \int K(y \mid \theta) dG_0(\theta)
+
\sum_{k} \frac{n_k}{\alpha+n} K(y \mid \theta_k^*)
\]

\subsection*{Interpretation}

\begin{itemize}
\item mixture of new component (prior) and existing clusters
\item adaptive smoothing depending on data
\item connects to kernel density estimation
\end{itemize}

\subsection*{Gibbs Sampling Strategy}

Update each $\theta_i$:
\[
p(\theta_i \mid \theta_{-i}, y)
\propto
\alpha p(y_i \mid \theta_i) G_0(\theta_i)
+
\sum_{j\neq i} p(y_i \mid \theta_j) \delta_{\theta_j}
\]

This forms basis for CRP-based samplers.

\subsection*{Key Insight}

DP mixtures provide:

\begin{itemize}
\item flexible density estimation
\item automatic clustering
\item uncertainty quantification on number of components
\end{itemize}

\section{Advanced MCMC Methods}

\subsection*{Metropolis--Hastings Algorithm}

Given current state $x^{(t)}$:

\[
Y_t \sim q(y \mid x^{(t)})
\]

Accept with probability
\[
\rho(x,y) = \min\left\{ \frac{\pi(y)\,q(x \mid y)}{\pi(x)\,q(y \mid x)}, 1 \right\}
\]

Update:
\[
x^{(t+1)} =
\begin{cases}
Y_t & \text{with prob } \rho(x^{(t)}, Y_t) \\
x^{(t)} & \text{otherwise}
\end{cases}
\]

\subsection*{Random Walk Metropolis--Hastings (RWMH)}

Proposal:
\[
Y_t = x^{(t)} + Z_t, \quad Z_t \sim g
\]

Thus
\[
q(y \mid x) = g(y - x)
\]

Common choices:
\[
Z_t \sim N(0, \sigma^2 I_d)
\]

\subsection*{Scaling of the Proposal}

Tradeoff:
\[
\sigma \text{ small} \Rightarrow \text{slow exploration}
\]
\[
\sigma \text{ large} \Rightarrow \text{high rejection}
\]

Goal: balance via acceptance rate.

\subsection*{Optimal Acceptance Rates}

For high dimension $d$:
\[
\text{Optimal rate} \approx 0.234
\]

For $d=1$:
\[
\text{Optimal rate} \approx 0.44
\]

\subsection*{Problems with MH}

\subsection*{Local Trap}

Example:
\[
\pi(x) = \frac{1}{3} N(\mu_1,\Sigma) + \frac{2}{3} N(\mu_2,\Sigma)
\]

Chains may remain in one mode:
\[
P(\text{jump between modes}) \approx \text{small}
\]

\subsection*{Intractable Normalizing Constants}

Target:
\[
f(x) \propto c(x)\psi(x)
\]

Acceptance ratio requires:
\[
\frac{c(y)}{c(x)}
\]

If $c(x)$ unknown $\Rightarrow$ MH not feasible.

\subsection*{Metropolis Adjusted Langevin Algorithm (MALA)}

Continuous-time:
\[
dX_t = dB_t + \frac{1}{2} \nabla \log \pi(X_t)
\]

Discrete proposal:
\[
Y_t \sim N\left(x^{(t)} + \frac{\sigma^2}{2} \nabla \log \pi(x^{(t)}), \sigma^2 I_d \right)
\]

\subsection*{Properties of MALA}

\[
\text{Optimal acceptance rate} \approx 0.574
\]

Advantages:
\[
\text{Uses gradient information} \Rightarrow \text{faster mixing}
\]

Limitation:
\[
\nabla \log \pi(x) \text{ must be computable}
\]

\subsection*{Implementation Issues}

Convergence diagnostics:
\[
\text{Need full exploration of support}
\]

Autocorrelation:
\[
\text{Low autocorrelation} \Rightarrow \text{efficient sampling}
\]

\subsection*{Multiple Chains vs Single Chain}

Multiple chains:
\[
\{x^{(m)}_t\}_{m=1}^M
\]

Single chain:
\[
x_1, x_2, \dots
\]

Goal:
\[
\text{independent samples} \approx \text{low autocorrelation}
\]

\subsection*{Thinning}

Keep every $k$-th sample:
\[
\{x_k, x_{2k}, \dots\}
\]

\subsection*{Reparametrization}

Problem:
\[
\text{High posterior correlation}
\]

Example:
\[
y_{ij} = \mu + \alpha_i + e_{ij}
\]

\[
\alpha_i \sim N(0, \tau^2), \quad e_{ij} \sim N(0,\sigma^2)
\]

Posterior correlations:
\[
\text{Cor}(\mu,\alpha_i) \approx -\left(1 + \frac{\sigma^2/n}{\tau^2/m}\right)^{-1/2}
\]

Reparametrize:
\[
\beta_i = \mu + \alpha_i
\]

Leads to reduced correlation and improved mixing.

\subsection*{Blocking}

Group parameters:
\[
\theta = (\theta_1,\dots,\theta_k)
\]

Instead of scalar updates:
\[
\theta_j \mid \theta_{-j}
\]

Use block updates:
\[
(\theta_{i_1}, \dots, \theta_{i_k}) \mid \theta_{-(i_1,\dots,i_k)}
\]

Benefit:
\[
\text{Better exploration in correlated directions}
\]

\subsection*{Collapsing}

Marginalize variables:
\[
p(x,y) \rightarrow p(x)
\]

Example:
\[
(X_1,X_2,X_3) \rightarrow (X_1,X_2)
\]

Then sample:
\[
X_3 \mid X_1,X_2
\]

Benefit:
\[
\text{Lower variance, faster convergence}
\]

\subsection*{Auxiliary Variable Methods}

Introduce $u$:
\[
p(x) = \int p(x,u)\,du
\]

Goal:
\[
\text{simplify sampling or avoid intractable terms}
\]

Examples:
\[
\text{Slice sampling}
\]

Key principle:
\[
\text{augment state space to simplify conditionals}
\]

\subsection*{Key Takeaways}

\[
\text{Efficiency depends on proposal design and parameterization}
\]

\[
\text{Use gradients (MALA), blocking, collapsing, or augmentation}
\]

\[
\text{Avoid local traps and high autocorrelation}
\]

\section{Auxiliary Variable Methods and Hybrid MCMC}

\subsection*{Cauchy--Normal Model via Data Augmentation}

\[
X_i \stackrel{iid}{\sim} \text{Cauchy}(\mu,1)
\]

Likelihood:
\[
L(\mu \mid x_{1:n}) = \prod_{i=1}^n \frac{1}{\pi(1 + (x_i - \mu)^2)}
\]

Prior:
\[
\mu \sim N(0,10)
\]

Posterior:
\[
\pi(\mu \mid x) \propto e^{-\mu^2/20} \prod_{i=1}^n \frac{1}{1 + (x_i - \mu)^2}
\]

Introduce latent variables:
\[
\frac{1}{1 + (x_i - \mu)^2}
=
\int_0^\infty e^{-\omega_i(1 + (x_i - \mu)^2)} d\omega_i
\]

Augmented posterior:
\[
\pi(\mu, \omega \mid x)
\propto
e^{-\mu^2/20}
\prod_{i=1}^n e^{-\omega_i(1 + (x_i - \mu)^2)}
\]

Gibbs updates:
\[
\omega_i \mid \mu \sim \text{Exp}(1 + (x_i - \mu)^2)
\]

\[
\mu \mid \omega \sim N\left(
\frac{\sum_i \omega_i x_i}{\sum_i \omega_i + 1/20},
\frac{1}{2\sum_i \omega_i + 1/10}
\right)
\]

\subsection*{Metropolis-within-Gibbs}

When a full conditional is not tractable:

\[
\theta_j \mid \theta_{-j} \quad \text{sampled via MH}
\]

Procedure:
\[
Y \sim q(\cdot \mid \theta_j)
\]

\[
\text{accept with prob } \min\left\{
\frac{\pi(Y \mid \theta_{-j}) q(\theta_j \mid Y)}
{\pi(\theta_j \mid \theta_{-j}) q(Y \mid \theta_j)}, 1
\right\}
\]

Key idea:
\[
\text{Replace difficult Gibbs steps with MH steps}
\]

\subsection*{Adaptive Gibbs Samplers}

Goal:
\[
\text{improve efficiency by learning during sampling}
\]

Selection probabilities:
\[
\alpha_n = (\alpha_{n,1}, \dots, \alpha_{n,d})
\]

Update rule:
\[
\alpha_n = R_n(\alpha_{n-1}, X_{n-1}, \dots)
\]

Transition:
\[
X_n \sim P_{\alpha_n}(X_{n-1}, \cdot)
\]

Key condition (diminishing adaptation):
\[
|\alpha_n - \alpha_{n-1}| \to 0
\]

Ensures:
\[
\text{ergodicity under suitable conditions}
\]

\subsection*{Random Scan Gibbs Sampler}

State:
\[
X = (X_1, \dots, X_d)
\]

At each step:
\[
i \sim \text{Categorical}(\alpha_1,\dots,\alpha_d)
\]

\[
X_i \sim \pi(\cdot \mid X_{-i})
\]

Kernel:
\[
P_\alpha = \sum_{i=1}^d \alpha_i P_i
\]

\subsection*{Adaptive Random Scan Gibbs}

Update:
\[
\alpha_n = R_n(\alpha_{n-1}, X_{0:n-1})
\]

Then:
\[
X_n \sim P_{\alpha_n}(X_{n-1}, \cdot)
\]

Key result:
\[
|\alpha_n - \alpha_{n-1}| \to 0
\quad \Rightarrow \quad \text{ergodicity (under conditions)}
\]

\subsection*{Metropolis-within-Gibbs (Adaptive)}

If conditional unavailable:
\[
X_i \sim Q_{x_{-i}}(X_i,\cdot)
\]

Accept:
\[
\min\left\{
1,
\frac{\pi(Y \mid x_{-i}) q(Y, X_i)}
{\pi(X_i \mid x_{-i}) q(X_i, Y)}
\right\}
\]

Adaptive version:
\[
\gamma_n \text{ controls proposal}
\]

\[
Q_{n} = Q_{\gamma_n}
\]

Condition:
\[
\|Q_{n+1} - Q_n\| \to 0
\]

\subsection*{Ergodicity Conditions (Key Structure)}

Two critical requirements:

Diminishing adaptation:
\[
\|P_n - P_{n-1}\| \to 0
\]

Containment:
\[
\text{chain does not escape to infinity}
\]

Together imply:
\[
X_n \xrightarrow{d} \pi
\]

\subsection*{Failure of Naive Adaptation}

Even if:
\[
\alpha_n \to \alpha
\]

Chain may still be:
\[
\text{non-ergodic (transient)}
\]

Hence:
\[
\text{adaptation must be controlled carefully}
\]

\subsection*{Adaptive Metropolis-within-Gibbs (Full)}

Updates both:
\[
\alpha_n \quad \text{and} \quad \gamma_n
\]

Algorithm:
\[
\alpha_n = R_n(\cdot), \quad \gamma_n = R'_n(\cdot)
\]

Then:
\[
X_n \sim P_{\alpha_n,\gamma_n}(X_{n-1}, \cdot)
\]

\subsection*{Slice Sampling (Revisited)}

Joint:
\[
p(x,u) = \frac{1}{Z} \mathbf{1}(0 < u < f(x))
\]

Gibbs steps:
\[
u \mid x \sim \text{Unif}(0, f(x))
\]

\[
x \mid u \sim \text{Unif}(\{x: f(x) \ge u\})
\]

For DP mixtures:
\[
u_i < w_{r_i}
\Rightarrow \text{finite truncation}
\]

\subsection*{Key Insight Across Methods}

All methods aim to:

\[
\text{improve mixing}
\]

via:
\[
\text{better proposals}, \quad \text{reparameterization}, \quad \text{augmentation}
\]

Unifying principle:
\[
\text{modify the Markov kernel while preserving } \pi
\]

\subsection*{Takeaways}

\[
\text{Gibbs: simple but may mix slowly}
\]

\[
\text{MH: flexible but sensitive to tuning}
\]

\[
\text{MALA: uses gradients}
\]

\[
\text{Adaptive MCMC: learns structure, but requires care}
\]

\[
\text{Auxiliary variables: turn hard problems into easy conditionals}
\]

\section{Pitman--Yor Process (PYP)}

\subsection*{Motivation: Limitations of the Dirichlet Process}

\begin{itemize}
\item DP mixture model:
\begin{itemize}
\item standard Bayesian nonparametric clustering model
\item exchangeable partitions
\item flexible number of clusters
\end{itemize}

\item Key limitations:
\begin{itemize}
\item cluster weights decay exponentially
\item number of clusters grows slowly:
\[
\mathbb{E}[K_n] \sim \alpha \log n
\]
\end{itemize}

\item Consequence:
\begin{itemize}
\item underestimates many small clusters
\item not suitable for heavy-tailed cluster structures
\end{itemize}
\end{itemize}

\subsection*{Idea of the Pitman--Yor Process}

\begin{itemize}
\item Goal:
\begin{itemize}
\item minimal extension of DP
\item preserve exchangeability
\item maintain tractable inference
\end{itemize}

\item Key modification:
\begin{itemize}
\item introduce discount parameter $d$
\item decouple:
\begin{itemize}
\item reinforcement (rich-get-richer)
\item tail behavior of cluster sizes
\end{itemize}
\end{itemize}

\item Result:
\begin{itemize}
\item exponential decay $\rightarrow$ power-law decay
\item heavier tails in cluster sizes
\end{itemize}
\end{itemize}

\subsection*{Definition: Stick-Breaking Representation}

\begin{itemize}
\item Random measure:
\[
G = \sum_{k=1}^\infty \pi_k \delta_{\phi_k}, \quad \phi_k \stackrel{iid}{\sim} G_0
\]

\item Weights:
\[
\pi_k = v_k \prod_{\ell < k} (1 - v_\ell)
\]

\item Stick-breaking variables:
\[
v_k \sim \text{Beta}(1 - d, \alpha + k d)
\]

\item Parameter constraints:
\[
0 \le d < 1, \quad \alpha > -d
\]

\item Special case:
\[
d = 0 \Rightarrow \text{Dirichlet Process}
\]
\end{itemize}

\subsection*{Pitman--Yor via Discounted CRP}

\begin{itemize}
\item Let:
\begin{itemize}
\item $K_n$ = number of clusters
\item $n_k$ = size of cluster $k$
\end{itemize}

\item Assignment probabilities:
\[
\mathbb{P}(z_{n+1} = k \mid z_{1:n}) =
\begin{cases}
\frac{n_k - d}{n + \alpha}, & k = 1, \dots, K_n \\
\frac{\alpha + d K_n}{n + \alpha}, & \text{new cluster}
\end{cases}
\]

\item Interpretation:
\begin{itemize}
\item reduced reinforcement: $n_k - d$ instead of $n_k$
\item increased innovation: $\alpha + d K_n$
\end{itemize}
\end{itemize}

\subsection*{Key Properties}

\begin{itemize}
\item Discounted reinforcement:
\begin{itemize}
\item weakens "rich-get-richer"
\item small clusters survive longer
\end{itemize}

\item Cluster-dependent innovation:
\begin{itemize}
\item more clusters $\Rightarrow$ higher probability of new cluster
\end{itemize}

\item Balance:
\begin{itemize}
\item prevents dominance of large clusters
\item encourages diversity
\end{itemize}
\end{itemize}

\subsection*{Power-Law Behavior of Weights}

\begin{itemize}
\item Expected weight:
\[
\mathbb{E}[\pi_k] = \mathbb{E}[v_k] \prod_{\ell<k} \mathbb{E}[1 - v_\ell]
\]

\item Moments:
\[
\mathbb{E}[v_k] = \frac{1 - d}{\alpha + kd + 1 - d}
\]

\[
\mathbb{E}[1 - v_\ell] = \frac{\alpha + \ell d}{\alpha + \ell d + 1 - d}
\]

\item Asymptotics:
\[
\mathbb{E}[\pi_k] \sim C_{\alpha,d} \, k^{-1/d}
\]

\item Interpretation:
\begin{itemize}
\item polynomial decay (heavy tail)
\item larger $d$ $\Rightarrow$ heavier tail
\end{itemize}
\end{itemize}

\subsection*{Growth of Number of Clusters}

\begin{itemize}
\item New cluster probability:
\[
\mathbb{P}(z_{n+1} = \text{new}) = \frac{\alpha + d K_n}{n + \alpha}
\]

\item Let $m_n = \mathbb{E}[K_n]$, then:
\[
m_{n+1} = m_n + \frac{\alpha + d m_n}{n + \alpha}
\]

\item Asymptotic approximation:
\[
\frac{dm_n}{dn} \approx \frac{d m_n}{n}
\]

\item Solve:
\[
\log m_n \approx d \log n
\quad \Rightarrow \quad
m_n \sim C n^d
\]

\item Compare:
\begin{itemize}
\item DP: $\mathbb{E}[K_n] \sim \log n$
\item PYP: $\mathbb{E}[K_n] \sim n^d$
\end{itemize}
\end{itemize}

\subsection*{Takeaways}

\begin{itemize}
\item PYP generalizes DP with one extra parameter $d$
\item introduces power-law cluster sizes
\item supports many small clusters
\item better fits real-world heavy-tailed data
\item retains tractable CRP-style inference
\end{itemize}

\chapter{Extensions of the Dirichlet Process}

\section{Pitman--Yor Process: Inference and Applications}

\subsection*{Inference under Pitman--Yor Process}

\begin{itemize}
\item Model:
\[
G \sim \text{PYP}(\alpha, d, G_0), \quad \theta_i \mid G \sim G, \quad y_i \mid \theta_i \sim F(\theta_i)
\]

\item Integrating out $G$ yields a collapsed Gibbs sampler similar to DP mixtures

\item Cluster assignment update:
\[
\mathbb{P}(z_i = k \mid z_{-i}, y) \propto
\begin{cases}
(n_{k,-i} - d)\, p(y_i \mid y_{k,-i}), & k \le K_{-i} \\
(\alpha + d K_{-i})\, p(y_i \mid G_0), & k = \text{new}
\end{cases}
\]

\item Same structure as DP mixture:
\begin{itemize}
\item only prior weights differ
\item likelihood terms unchanged
\end{itemize}

\item Computational cost:
\begin{itemize}
\item comparable to DP inference
\end{itemize}
\end{itemize}

\subsection*{DP vs PYP: Key Differences}

\begin{itemize}
\item Parameters:
\[
\text{DP: } \alpha > 0, \quad \text{PYP: } \alpha > -d,\; 0 \le d < 1
\]

\item Stick-breaking:
\[
\text{DP: } v_k \sim \text{Beta}(1,\alpha), \quad
\text{PYP: } v_k \sim \text{Beta}(1-d,\alpha+kd)
\]

\item Weight decay:
\[
\text{DP: exponential}, \quad \text{PYP: power-law}
\]

\item Cluster growth:
\[
\mathbb{E}[K_n] \sim \log n \quad \text{(DP)}, \qquad \mathbb{E}[K_n] \sim n^d \quad \text{(PYP)}
\]

\item Reinforcement:
\[
\text{DP: } n_k, \quad \text{PYP: } n_k - d
\]

\item New cluster probability:
\[
\text{DP: } \frac{\alpha}{n+\alpha}, \quad
\text{PYP: } \frac{\alpha + dK_n}{n+\alpha}
\]
\end{itemize}

\subsection*{When to Use Pitman--Yor}

\begin{itemize}
\item DP assumptions:
\begin{itemize}
\item new clusters quickly become unlikely
\item tail mass decays exponentially
\item small clusters disappear
\end{itemize}

\item Use PYP when:
\begin{itemize}
\item persistent rare categories
\item heavy-tailed cluster sizes
\item continual discovery of new types
\end{itemize}

\item Applications:
\begin{itemize}
\item genetics (rare alleles)
\item text modeling (long-tail words)
\item ecology (species abundance)
\item social networks (degree distributions)
\end{itemize}
\end{itemize}

\subsection*{Case Study: Genetics}

\begin{itemize}
\item Setup:
\begin{itemize}
\item $n$ individuals, each with allele at a locus
\item $K_n$ distinct alleles, counts $n_k$
\end{itemize}

\item Empirical pattern:
\begin{itemize}
\item few common alleles
\item many rare alleles
\end{itemize}

\item Model:
\[
G \sim \text{PYP}(\alpha, d, G_0), \quad z_i \mid G \sim G
\]

\item Benefit:
\begin{itemize}
\item heavy-tailed cluster sizes preserve rare alleles
\end{itemize}
\end{itemize}

\subsection*{Case Study: Topic Modeling}

\begin{itemize}
\item Topic distributions:
\[
\phi_k = (\phi_{k,1}, \phi_{k,2}, \dots), \quad \sum_v \phi_{k,v} = 1
\]

\item Word generation:
\[
w_i \mid z_i = k \sim \text{Categorical}(\phi_k)
\]

\item Empirical observation:
\begin{itemize}
\item few very frequent words
\item many rare words
\end{itemize}

\item Model:
\[
\phi_k \sim \text{PYP}(\alpha, d, G_0)
\]

\item Effect:
\begin{itemize}
\item heavy-tailed word distributions
\item preserves rare words
\end{itemize}
\end{itemize}

\section{Hierarchical Dirichlet Process (HDP)}

\subsection*{Motivation: Multiple Related Groups}

\begin{itemize}
\item Data partitioned into groups:
\[
\{x_{ji}\}, \quad j=1,\dots,J,\; i=1,\dots,n_j
\]

\item Within each group:
\begin{itemize}
\item observations are exchangeable
\item each group has its own mixture distribution
\end{itemize}

\item Goal:
\begin{itemize}
\item unknown number of components
\item components shared across groups
\item group-specific weights
\end{itemize}
\end{itemize}

\subsection*{Why Sharing Matters}

\begin{itemize}
\item Unknown number of clusters:
\begin{itemize}
\item parametric models require fixing $K$
\item nonparametric models allow growth with data
\end{itemize}

\item Shared components:
\begin{itemize}
\item groups are not independent
\item large groups inform small ones
\end{itemize}

\item Different weights:
\begin{itemize}
\item prevalence varies across groups
\end{itemize}
\end{itemize}

\subsection*{Failure of Naive Construction}

\begin{itemize}
\item Attempt:
\[
G_j \sim \text{DP}(\alpha_j, G_0(\tau)), \quad \tau \sim p(\tau)
\]

\item Conditional independence:
\[
G_j \perp G_{j'} \mid \tau
\]

\item Problem:
\begin{itemize}
\item if $G_0(\tau)$ is continuous:
\[
\mathbb{P}(\theta_{jk} = \theta_{j'k'}) = 0
\]
\item no shared atoms
\end{itemize}

\item Conclusion:
\begin{itemize}
\item sharing parameter $\tau$ is insufficient
\item need discrete base measure
\end{itemize}
\end{itemize}

\subsection*{HDP Model Definition}

\begin{itemize}
\item Global measure:
\[
G_0 \mid \gamma, H \sim \text{DP}(\gamma, H)
\]

\item Group-specific measures:
\[
G_j \mid \alpha_0, G_0 \sim \text{DP}(\alpha_0, G_0)
\]

\item Observations:
\[
\theta_{ji} \mid G_j \sim G_j, \quad x_{ji} \mid \theta_{ji} \sim F(\theta_{ji})
\]

\item Key property:
\begin{itemize}
\item $G_0$ is discrete
\item atoms are shared across all $G_j$
\end{itemize}
\end{itemize}

\section{Hierarchical Dirichlet Process: Constructions and Inference}

\subsection*{Graphical Interpretation and Limitation of Naive Model}

\begin{itemize}
\item Naive construction:
\[
G_j \sim \text{DP}(\alpha_j, G_0(\tau)), \quad \tau \sim p(\tau)
\]

\item Conditional independence:
\[
G_j \perp G_{j'} \mid \tau
\]

\item Even with shared hyperparameter $\tau$, atoms differ across groups:
\begin{itemize}
\item $G_j$ are independent draws
\item no shared support
\end{itemize}

\item Conclusion:
\begin{itemize}
\item no sharing of mixture components
\item motivates hierarchical construction
\end{itemize}
\end{itemize}

\subsection*{Stick-Breaking Construction of HDP}

\begin{itemize}
\item Global stick-breaking:
\[
G_0 = \sum_{k=1}^{\infty} \beta_k \delta_{\phi_k}, \quad \phi_k \overset{iid}{\sim} H, \quad \beta \sim \text{GEM}(\gamma)
\]

\item Group-level measures:
\[
G_j \mid G_0 \sim \text{DP}(\alpha_0, G_0)
\]

\item Since $G_0$ is discrete, all $G_j$ share atoms $\{\phi_k\}$:
\[
G_j = \sum_{k=1}^{\infty} \pi_{jk} \delta_{\phi_k}
\]

\item Goal:
\begin{itemize}
\item derive stick-breaking form for $\pi_{jk}$
\end{itemize}
\end{itemize}

\subsection*{Finite-Dimensional Marginals}

\begin{itemize}
\item For partition $(A_1,\dots,A_r)$:
\[
(G_j(A_1),\dots,G_j(A_r)) \sim \text{Dir}(\alpha_0 G_0(A_1),\dots,\alpha_0 G_0(A_r))
\]

\item Using decomposition:
\[
G_j(A_\ell) = \sum_{k \in K_\ell} \pi_{jk}, \quad G_0(A_\ell) = \sum_{k \in K_\ell} \beta_k
\]

\item Leads to:
\[
\left(\sum_{k \in K_1} \pi_{jk}, \dots, \sum_{k \in K_r} \pi_{jk} \right)
\sim \text{Dir}\left(\alpha_0 \sum_{k \in K_1} \beta_k, \dots \right)
\]
\end{itemize}

\subsection*{Stick-Breaking for Group Weights}

\begin{itemize}
\item Representation:
\[
\pi_{jk} = \pi'_{jk} \prod_{\ell<k} (1 - \pi'_{j\ell})
\]

\item Where:
\[
\pi'_{jk} \sim \text{Beta}\left(\alpha_0 \beta_k,\; \alpha_0 \sum_{\ell>k} \beta_\ell \right)
\]

\item Derived using:
\begin{itemize}
\item Dirichlet aggregation
\item renormalization properties
\end{itemize}
\end{itemize}

\subsection*{Chinese Restaurant Franchise (CRF)}

\begin{itemize}
\item Hierarchical model:
\[
G_0 \sim \text{DP}(\gamma, H), \quad G_j \sim \text{DP}(\alpha_0, G_0)
\]

\item Interpretation:
\begin{itemize}
\item each group = restaurant
\item observations = customers
\item tables = local clusters
\item dishes = global clusters
\end{itemize}
\end{itemize}

\subsection*{CRF Variables}

\begin{itemize}
\item Global dishes:
\[
\phi_k \sim H
\]

\item Table-level assignments:
\[
\psi_{jt} = \phi_{k_{jt}}
\]

\item Customer-level:
\[
\theta_{ji} = \psi_{j,t_{ji}}
\]

\item Observations:
\[
x_{ji} \sim F(\theta_{ji})
\]
\end{itemize}

\subsection*{Predictive Distribution for Customers}

\begin{itemize}
\item Integrating out $G_j$:
\[
\theta_{ji} \mid \theta_{j,1:i-1} \sim \sum_{t=1}^{m_j} \frac{n_{jt}}{i-1+\alpha_0}\delta_{\psi_{jt}} + \frac{\alpha_0}{i-1+\alpha_0} G_0
\]

\item Interpretation:
\begin{itemize}
\item join existing table with prob $\propto n_{jt}$
\item start new table with prob $\propto \alpha_0$
\end{itemize}
\end{itemize}

\subsection*{Predictive Distribution for Tables}

\begin{itemize}
\item Integrating out $G_0$:
\[
\psi_{jt} \mid \{\psi\}_{-jt} \sim \sum_{k=1}^{K} \frac{m_k}{m_{\cdot}+\gamma}\delta_{\phi_k} + \frac{\gamma}{m_{\cdot}+\gamma} H
\]

\item Where:
\begin{itemize}
\item $m_k$ = number of tables serving dish $k$
\item $m_{\cdot} = \sum_k m_k$
\end{itemize}

\item Interpretation:
\begin{itemize}
\item reuse global dish with prob $\propto m_k$
\item create new dish with prob $\propto \gamma$
\end{itemize}
\end{itemize}

\subsection*{Index-Based Sampling View}

\begin{itemize}
\item Customer $\rightarrow$ table:
\[
\mathbb{P}(t_{ji} = t \mid \cdot) =
\begin{cases}
\frac{n_{jt}}{i-1+\alpha_0}, & t \le m_j \\
\frac{\alpha_0}{i-1+\alpha_0}, & t = m_j+1
\end{cases}
\]

\item Table $\rightarrow$ dish:
\[
\mathbb{P}(k_{jt} = k \mid \cdot) =
\begin{cases}
\frac{m_k}{m_{\cdot}+\gamma}, & k \le K \\
\frac{\gamma}{m_{\cdot}+\gamma}, & k = K+1
\end{cases}
\]
\end{itemize}

\subsection*{Collapsed Gibbs Sampling: Table Assignment Update}

\begin{itemize}
\item For customer $i$ in group $j$, the table assignment update has the Gibbs form
\[
p(t_{ji}=t \mid \text{rest}, x)
\propto
p(t_{ji}=t \mid \text{rest}) \,
p(x_{ji} \mid t_{ji}=t,\text{rest})
\]

\item The prior term comes from the CRP within group $j$:
\[
p(t_{ji}=t \mid t^{-ji})
\propto
\begin{cases}
n_{jt}^{-ji}, & \text{if table } t \text{ exists}, \\
\alpha_0, & \text{if table } t \text{ is new}.
\end{cases}
\]

\item The likelihood term is collapsed by integrating out the dish parameter $\phi_k$.
\end{itemize}

\subsection*{Collapsed Predictive Density}

Define the collapsed predictive density
\[
f_k^{-x_{ji}}(x_{ji})
=
p(x_{ji} \mid \{x_{j'i'} : z_{j'i'}=k,\; (j',i') \neq (j,i)\}).
\]

Using Bayes' rule,
\[
f_k^{-x_{ji}}(x_{ji})
=
\frac{
\int f(x_{ji}\mid \phi_k)
\prod_{j'i'\neq ji,\; z_{j'i'}=k}
f(x_{j'i'}\mid \phi_k)
h(\phi_k)\,d\phi_k
}{
\int
\prod_{j'i'\neq ji,\; z_{j'i'}=k}
f(x_{j'i'}\mid \phi_k)
h(\phi_k)\,d\phi_k
}.
\]

\begin{itemize}
\item $h$ is the density of the base measure $H$.

\item The denominator is the marginal likelihood of observations currently assigned to global cluster $k$, excluding $x_{ji}$.

\item The numerator is the marginal likelihood if $x_{ji}$ were added to cluster $k$.

\item If $h(\phi)$ is conjugate to $f(x\mid \phi)$, this predictive density is available in closed form.
\end{itemize}

\subsection*{Collapsed Gibbs Sampling: Dish Assignment Update}

\begin{itemize}
\item For a newly created table $t$ in group $j$, the dish assignment update is
\[
p(k_{jt}=k \mid \text{rest}, x)
\propto
p(k_{jt}=k \mid \text{rest})\,
p(x_{jt} \mid k_{jt}=k,\text{rest}).
\]

\item The prior term comes from the global CRP:
\[
p(k_{jt}=k \mid k^{-jt})
\propto
\begin{cases}
m_{\cdot k}, & \text{if dish } k \text{ exists}, \\
\gamma, & \text{if dish } k \text{ is new}.
\end{cases}
\]

\item The likelihood term is collapsed over $\phi_k$. If table $t$ is assigned dish $k$, all customers at that table share parameter $\phi_k$:
\[
p(x_{jt}\mid k,\text{rest})
=
\prod_{i:t_{ji}=t}
f_k^{-x_{ji}}(x_{ji}).
\]

\item Since the table was just created, it contains the single customer that triggered the new table:
\[
\prod_{i:t_{ji}=t} f_k^{-x_{ji}}(x_{ji})
=
f_k^{-x_{ji}}(x_{ji}).
\]

\item Therefore:
\[
p(k_{jt}=k \mid \cdot)
\propto
\begin{cases}
m_{\cdot k}\, f_k^{-x_{ji}}(x_{ji}), & \text{existing dish } k, \\
\gamma\, f_{\text{new}}^{-x_{ji}}(x_{ji}), & \text{new dish}.
\end{cases}
\]
\end{itemize}

\subsection*{Case Study: BNP Modeling of Marathon Data}

\begin{itemize}
\item Context:
\begin{itemize}
\item Marathon finishing times depend on age, gender, and race conditions.
\item Traditional grading methods rely on top records and coarse adjustments.
\item Goal: data-driven comparison using latent structure.
\end{itemize}

\item Problem:
\begin{itemize}
\item Discover latent \emph{running patterns} (pace profiles across race segments).
\item Cluster runners within each age/gender group.
\item Share patterns across groups using an HDP.
\end{itemize}

\item Setup:
\begin{itemize}
\item Groups indexed by $j$ (age/gender).
\item Runner $i$ in group $j$ represented by:
\[
x_{ji} = (\text{fraction of time spent in each segment}).
\]
\end{itemize}

\item HDP model:
\begin{itemize}
\item Global distribution:
\[
G_0 \sim \text{DP}(\gamma, H)
\]

\item Group-level:
\[
G_j \sim \text{DP}(\alpha, G_0), \quad
G_j = \sum_k \pi_{jk} \delta_{\phi_k}
\]

\item Likelihood:
\[
x_{ji} \mid z_{ji}=k \sim \text{Dirichlet}(\tau \phi_k)
\]
\end{itemize}

\item Interpretation:
\begin{itemize}
\item $\phi_k$ = global pace pattern
\item $\pi_{jk}$ = prevalence of pattern in group $j$
\end{itemize}

\item Data:
\begin{itemize}
\item NYC Marathon (2007--2011)
\item $\sim 194{,}000$ runners
\end{itemize}

\item Outcome:
\begin{itemize}
\item $\sim 46$ latent pace clusters discovered
\item Examples:
\begin{itemize}
\item fast-start then slow-down
\item steady pace
\end{itemize}
\end{itemize}

\item Insight:
\begin{itemize}
\item Groups differ in mixture weights but share global patterns
\item Enables prediction of finishing time from partial race data
\end{itemize}
\end{itemize}

\section{Dependent Dirichlet Processes (DDP)}

\subsection*{Motivation}

\begin{itemize}
\item Standard DP assumes a single distribution:
\[
G \sim \text{DP}(\alpha, H)
\]

\item In regression:
\[
y_i = f(x_i) + \epsilon_i
\]

\item Need distributions that vary with $x$:
\[
G_x \quad \text{depends on } x
\]
\end{itemize}

\subsection*{DDP Construction}

\begin{itemize}
\item Define a collection:
\[
\mathcal{G} = \{G_x : x \in \mathcal{X}\}
\]

\item Marginals:
\[
G_x \sim \text{DP}(\cdot,\cdot)
\]

\item Stick-breaking representation:
\[
G_x = \sum_h p_h \delta_{m_h(x)}
\]

\item Key idea:
\begin{itemize}
\item weights $p_h$ may be shared
\item atoms $m_h(x)$ vary with $x$
\end{itemize}
\end{itemize}

\subsection*{Types of Dependence}

\begin{itemize}
\item Common weights, dependent atoms:
\[
p_h \text{ fixed}, \quad m_h(x) \text{ varies with } x
\]

\item Common atoms, dependent weights

\item Both dependent

\item Partial sharing via discrete base measure
\end{itemize}

\subsection*{Example: Health Outcomes}

\begin{itemize}
\item Variables:
\begin{itemize}
\item $Y$ = health outcome
\item $x_1$ = exposure
\item $x_2,\dots,x_p$ = covariates
\end{itemize}

\item Interest:
\[
\text{distribution of } Y \mid x
\]

\item Motivation:
\begin{itemize}
\item effects may appear in tails
\item avoid dichotomizing outcomes
\end{itemize}
\end{itemize}

\subsection*{DDE and Preterm Birth Example}

\begin{itemize}
\item Logistic model:
\[
\log \Pr(y_i=1 \mid x_i,\beta)
=
\beta_0 + \sum_{h=1}^4 \beta_h\, dde_{hi} + z_i^\top \alpha
\]

\item Interpretation:
\begin{itemize}
\item exposure categories modeled via indicators $dde_{hi}$
\item $z_i$ = confounders
\end{itemize}

\item Results:
\begin{itemize}
\item increasing odds ratios across exposure levels
\item strong evidence of dose-response relationship
\end{itemize}
\end{itemize}

\subsection*{Key Insight of DDP}

\begin{itemize}
\item Instead of modeling only the mean:
\[
\mathbb{E}[Y \mid x]
\]

\item Model full conditional distribution:
\[
Y \mid x \sim G_x
\]

\item Allows:
\begin{itemize}
\item heteroskedasticity
\item multimodality
\item tail effects
\end{itemize}
\end{itemize}

\subsection{Mixture Models and Regression Extensions}

\subsubsection*{Basic Mixture Model}
\begin{itemize}
\item Let $y_i$ denote gestational age for individual $i$.
\item Model density as a mixture:
\[
f(y_i) = \int \mathcal{N}(y_i; \mu, \sigma^2)\, dG(\mu, \sigma^2)
\]
\item Finite mixture approximation:
\[
G = \sum_{h=1}^k p_h \delta_{\theta_h}, \quad \theta_h = (\mu_h, \sigma_h^2)
\]
\item Mixtures of Gaussians can approximate flexible densities.
\end{itemize}

\subsubsection*{Mixtures of Regression Models}
\begin{itemize}
\item Incorporate predictors $x_i$ (e.g., DDE):
\[
f(y_i \mid x_i) = \sum_{h=1}^k p_h \mathcal{N}(y_i; x_i'\beta_h, \sigma_h^2)
\]
\item Hierarchical mixture of experts (HME):
\begin{itemize}
\item Replace weights $p_h$ with $p_h(x_i)$
\item Allows predictor-dependent clustering
\end{itemize}
\item Typically estimated via EM:
\begin{itemize}
\item Can overfit
\item Requires choosing $k$
\end{itemize}
\end{itemize}

\subsubsection*{Infinite Mixtures Motivation}
\begin{itemize}
\item Avoid fixing number of components $k$
\item Introduce nonparametric prior:
\[
G_x = \{G_x : x \in \mathcal{X}\}
\]
\item If $G$ is global, use:
\[
G \sim \mathrm{DP}(\alpha G_0)
\]
\end{itemize}

\subsection{Dependent Dirichlet Processes (DDP)}

\subsubsection*{General Form}
\begin{itemize}
\item Define predictor-dependent random measure:
\[
G_x = \sum_{h=1}^\infty p_h \delta_{\theta_h(x)}
\]
\item Atoms vary with $x$:
\[
\theta_h(\cdot) \sim \text{Gaussian Process}
\]
\item Weights may be:
\begin{itemize}
\item shared across $x$
\item or predictor-dependent
\end{itemize}
\end{itemize}

\subsubsection*{Key Idea}
\begin{itemize}
\item Introduce dependence across $x$ through:
\begin{itemize}
\item weights $p_h(x)$, or
\item locations $\theta_h(x)$
\end{itemize}
\item Goal: flexible conditional distributions $f(y \mid x)$
\end{itemize}

\subsection{Weighted Mixture of Dirichlet Processes}

\begin{itemize}
\item Construct local DPs at knots $x_j^*$:
\[
G_j^* \sim \mathrm{DP}(M, G_0)
\]
\item Define:
\[
G_x(B) = \sum_{j=1}^J \frac{K(x, x_j^*)}{\sum_\ell K(x, x_\ell^*)} G_j^*(B)
\]
\item Interpretation:
\begin{itemize}
\item Nearby knots receive higher weight
\item Smooth variation across predictor space
\end{itemize}
\end{itemize}

\subsection{Kernel Stick-Breaking Process (KSBP)}

\subsubsection*{Construction}
\begin{itemize}
\item Define:
\[
G_x = \sum_{h=1}^\infty \left[ V_h K(x, \Gamma_h) \prod_{\ell<h}(1 - V_\ell K(x, \Gamma_\ell)) \right] G_h^*
\]
\item Components:
\begin{itemize}
\item $V_h \sim \text{Beta}(a_h, b_h)$
\item $\Gamma_h \sim H$ (locations)
\item $G_h^* \sim \mathrm{DP}(\alpha G_0)$
\end{itemize}
\item $K(x,\Gamma_h)$ is a bounded kernel
\end{itemize}

\subsubsection*{Interpretation}
\begin{itemize}
\item Weights depend on distance between $x$ and $\Gamma_h$
\item Nearby components get larger weights
\item Induces smooth dependence across predictors
\end{itemize}

\subsubsection*{Special Cases}
\begin{itemize}
\item If $K(x,\Gamma_h)=1$:
\begin{itemize}
\item recovers standard DP or Pitman--Yor (depending on $V_h$)
\end{itemize}
\item If $K(x,\Gamma_h)=\exp(-\psi\|x-\Gamma_h\|)$:
\begin{itemize}
\item weights decay with distance
\end{itemize}
\end{itemize}

\subsection{Predictor-Dependent Clustering}

\begin{itemize}
\item Marginal predictive:
\[
(\phi_i \mid \phi^{(i)}, X)
=
\frac{\alpha}{\alpha + w_i} G_0
+
\sum_{j \neq i} \frac{w_{ij}}{\alpha + w_i} \delta_{\phi_j}
\]
\item Weights $w_{ij}$ depend on similarity in predictors
\item Key difference from DP:
\begin{itemize}
\item clustering depends on $x$
\end{itemize}
\end{itemize}

\subsection{Hierarchical Dirichlet Process (HDP) Revisited}

\begin{itemize}
\item Model:
\[
G_0 \sim \mathrm{DP}(\gamma, H), \quad G_j \mid G_0 \sim \mathrm{DP}(\alpha, G_0)
\]
\item Interpretation:
\begin{itemize}
\item $G_0$ defines global atoms
\item $G_j$ selects subsets with group-specific weights
\end{itemize}
\item Enables sharing across groups without explicit regression
\end{itemize}

\subsection{Example: Information Retrieval}

\begin{itemize}
\item Words:
\[
x_i \sim G_j
\]
\item $G_j$: document-specific distribution
\item $G_0$: global vocabulary distribution
\item Topics shared across documents
\end{itemize}

\subsection{Nested Dirichlet Process}

\begin{itemize}
\item Model:
\begin{itemize}
\item $G_j$ themselves are clustered
\item $G_j \sim Q$, where $Q$ is a DP over distributions
\end{itemize}
\item Two levels:
\begin{itemize}
\item clustering of groups (e.g., states)
\item clustering within groups (e.g., hospitals)
\end{itemize}
\item Allows discovery of higher-level structure
\end{itemize}

\subsection{ANOVA Dependent Dirichlet Process (ANOVA-DDP)}

\subsubsection*{Model Construction}
\begin{itemize}
\item Consider categorical predictors:
\[
x_j = (v, w)
\]
(e.g., row/column effects in ANOVA)

\item Introduce dependence across $x$ via ANOVA structure:
\[
m_{xh} = \mu_h + \alpha_{h,v} + \beta_{h,w}
\]

\item Each mixture component varies with factor levels while sharing global structure.
\end{itemize}

\subsubsection*{Equivalent Representation}
\begin{itemize}
\item Random measure:
\[
G_x(\theta) = \sum_h w_h \delta_{m_{xh}}
\]

\item Can be written as:
\[
\theta_{xi} = \Phi_i' d_x
\]

\item Where:
\begin{itemize}
\item $d_x$ = design vector (ANOVA encoding)
\item $M_h = (\mu_h, \alpha_{1,h}, \ldots, \alpha_{V,h}, \beta_{1,h}, \ldots, \beta_{W,h})$
\end{itemize}

\item Interpretation:
\begin{itemize}
\item ANOVA-DDP = DP mixture of ANOVA models
\end{itemize}
\end{itemize}

\subsection{Problem Setting for BNP Regression}

\begin{itemize}
\item Data:
\[
\{(y_i, x_i)\}_{i=1}^n, \quad x_i \in \mathbb{R}^Q
\]

\item Model:
\[
y_i = f(x_i) + \epsilon_i, \quad \epsilon_i \sim \mathcal{N}(0,\sigma^2)
\]

\item Objective:
\begin{itemize}
\item Learn a distribution over functions $p(f)$
\end{itemize}
\end{itemize}

\subsubsection*{Applications}
\begin{itemize}
\item Regression (curve fitting, smoothing)
\item Classification (categorical outputs)
\item Prediction (generalization)
\end{itemize}

\subsection{Gaussian Process (GP)}

\subsubsection*{Definition}
\begin{itemize}
\item A Gaussian process is a distribution over functions:
\[
f \sim \mathrm{GP}(m(x), K(x,x'))
\]

\item Any finite collection $\{f(x_1), \ldots, f(x_n)\}$ is jointly Gaussian.
\end{itemize}

\subsubsection*{Covariance Structure}
\begin{itemize}
\item Covariance matrix:
\[
K_{ij} = K(x_i, x_j)
\]

\item Interpretation:
\[
K(x_i, x_j) = \mathbb{E}[f(x_i) f(x_j)]
\]

\item Must be positive semidefinite.
\end{itemize}

\subsection{Common Covariance Functions}

\subsubsection*{Squared Exponential (SE)}
\begin{itemize}
\item Kernel:
\[
K(x,x') = \sigma_0^2 \exp\left(-\frac{1}{2}\left(\frac{x-x'}{\lambda}\right)^2\right)
\]

\item Properties:
\begin{itemize}
\item Smooth (infinitely differentiable)
\item Stationary
\end{itemize}

\item Hyperparameters:
\begin{itemize}
\item $\lambda$: lengthscale
\item $\sigma_0$: amplitude
\end{itemize}
\end{itemize}

\subsubsection*{Matérn Class}
\begin{itemize}
\item Kernel:
\[
K(x,x') =
\frac{2^{1-\nu}}{\Gamma(\nu)}
\left(\frac{\sqrt{2\nu}|x-x'|}{\lambda}\right)^\nu
K_\nu\left(\frac{\sqrt{2\nu}|x-x'|}{\lambda}\right)
\]

\item Controls smoothness:
\begin{itemize}
\item $\nu \to \infty$: SE kernel
\item $\nu = 1/2$: exponential kernel
\end{itemize}
\end{itemize}

\subsubsection*{Other Kernels}
\begin{itemize}
\item Linear:
\[
K(x,x') = \sigma_0^2 + x x'
\]

\item Brownian motion:
\[
K(x,x') = \min(x,x')
\]

\item Periodic:
\[
K(x,x') = \exp\left(-\frac{2\sin^2((x-x')/2)}{\lambda^2}\right)
\]
\end{itemize}

\subsection{Combining Kernels}

\begin{itemize}
\item Sum:
\[
K = K_1 + K_2
\]

\item Product:
\[
K = K_1 K_2
\]

\item Convolution:
\[
K(x,x') = \int h(x,z) K(z,z') h(x',z')\, dz\, dz'
\]
\end{itemize}

\subsection{GP Regression}

\begin{itemize}
\item Prior:
\[
f \sim \mathrm{GP}(m, K)
\]

\item Likelihood:
\[
y \mid f \sim \mathcal{N}(f(X), \sigma^2 I)
\]

\item Posterior:
\[
f \mid y \sim \mathrm{GP}(m_p, K_p)
\]

\item Posterior mean:
\[
m_p(x) = m(x) + K(X,x)^T (K + \sigma^2 I)^{-1}(y - m(X))
\]

\item Posterior covariance:
\[
K_p(x,x') = K(x,x') - K(X,x)^T (K + \sigma^2 I)^{-1} K(X,x')
\]
\end{itemize}

\subsection{Prediction (Multi-step)}

\begin{itemize}
\item Joint Gaussian:
\[
\begin{bmatrix} y \\ y^* \end{bmatrix}
\sim \mathcal{N}
\left(
\begin{bmatrix} \mu_1 \\ \mu_2 \end{bmatrix},
\begin{bmatrix} \Sigma_{11} & \Sigma_{12} \\ \Sigma_{21} & \Sigma_{22} \end{bmatrix}
\right)
\]

\item Conditional mean:
\[
\mathbb{E}[y^* \mid y] = \mu_2 + \Sigma_{21}\Sigma_{11}^{-1}(y - \mu_1)
\]

\item Conditional variance:
\[
\mathrm{Var}(y^* \mid y) = \Sigma_{22} - \Sigma_{21}\Sigma_{11}^{-1}\Sigma_{12}
\]
\end{itemize}

\subsection{Learning Hyperparameters}

\begin{itemize}
\item Marginal likelihood:
\[
p(y \mid X, \theta) = \mathcal{N}(0, K_\theta + \sigma^2 I)
\]

\item Log marginal likelihood:
\[
\log p(y \mid X, \theta)
=
-\frac{1}{2}\log |K_\theta + \sigma^2 I|
-\frac{1}{2} y^T (K_\theta + \sigma^2 I)^{-1} y
+ \text{const}
\]

\item Estimation:
\begin{itemize}
\item Maximum likelihood
\item Fully Bayesian (e.g., HMC)
\end{itemize}
\end{itemize}

\subsection{Remarks}

\begin{itemize}
\item GPs:
\begin{itemize}
\item handle uncertainty via averaging
\item learn kernel structure from data
\item avoid manual feature engineering
\end{itemize}

\item Compared to SVM:
\begin{itemize}
\item probabilistic (uncertainty quantification)
\item automatic regularization
\end{itemize}
\end{itemize}

\subsection{Feature Selection}

\begin{itemize}
\item Model selection problem:
\[
m \in \{0,1\}^D
\]

\item Ideal solution:
\[
\hat{m} = \arg\max_m p(D \mid m)
\]

\item Challenges:
\begin{itemize}
\item combinatorial ($2^D$ models)
\item overfitting
\end{itemize}

\item Bayesian solution:
\begin{itemize}
\item average over models instead of selecting one
\end{itemize}
\end{itemize}

\subsection{Feature Selection in Bayesian Models}

\begin{itemize}
\item Key questions:
\begin{itemize}
\item Why perform feature selection?
\item What is the cost of including all features?
\end{itemize}

\item Important distinction:
\begin{itemize}
\item Overfitting is less relevant in fully Bayesian methods (no point estimation)
\item Bayesian averaging already regularizes models
\end{itemize}

\item Feature selection is justified when:
\begin{itemize}
\item Measurement cost is high
\item Prediction with many features is computationally expensive
\end{itemize}

\item Empirical note:
\begin{itemize}
\item Bayesian methods can perform well using all features
\end{itemize}
\end{itemize}

\subsubsection*{Automatic Relevance Determination (ARD)}

\begin{itemize}
\item Kernel with dimension-specific lengthscales:
\[
K(x_n, x_{n'}) = \nu \exp\left(-\frac{1}{2} \sum_{d=1}^D \left(\frac{x_n^{(d)} - x_{n'}^{(d)}}{r_d}\right)^2\right)
\]

\item Interpretation:
\begin{itemize}
\item $r_d$ controls variation along dimension $d$
\item $r_d \to \infty$ $\Rightarrow$ dimension becomes irrelevant
\end{itemize}

\item Learning $\{r_d\}_{d=1}^D$ enables automatic feature selection.
\end{itemize}

\subsection{Limitations of Gaussian Processes}

\begin{itemize}
\item Computational cost:
\[
\mathcal{O}(n^3)
\]
due to matrix inversion

\item Poor handling of:
\begin{itemize}
\item discontinuities
\item sharp structural changes
\end{itemize}

\item Smooth kernels (e.g., RBF) may be overly restrictive.
\end{itemize}

\subsection{Scalable Gaussian Processes}

\subsubsection*{Global Approximations}

\begin{itemize}
\item Use subset of size $m \ll n$:
\begin{itemize}
\item random subsampling
\item clustering (e.g., centroids)
\item active learning selection
\end{itemize}

\item Low-rank approximation:
\[
K_{nn} \approx K_{nm} K_{mm}^{-1} K_{mn}
\]
\end{itemize}

\subsubsection*{Sparse Kernels}

\begin{itemize}
\item Compact support kernels:
\[
K(x_i, x_j) = 0 \quad \text{if } \|x_i - x_j\| > c
\]

\item Challenges:
\begin{itemize}
\item maintain positive semidefiniteness
\end{itemize}

\item Complexity:
\[
\mathcal{O}(\alpha n^3)
\]
with sparsity factor $\alpha$.
\end{itemize}

\subsubsection*{Sparse GP (Inducing Points)}

\begin{itemize}
\item Introduce $M \ll N$ inducing points

\item Approximate prior:
\[
p(f) \approx \mathcal{N}(0, K_{NM} K_{MM}^{-1} K_{MN} + \Lambda)
\]

\item Benefits:
\begin{itemize}
\item Faster inversion:
\[
\mathcal{O}(M^2 N) \ll \mathcal{O}(N^3)
\]
\item Efficient prediction:
\[
\mathcal{O}(M) \text{ mean}, \quad \mathcal{O}(M^2) \text{ variance}
\]
\end{itemize}
\end{itemize}

\subsection{CART: Regression Trees}

\begin{itemize}
\item Motivation:
\begin{itemize}
\item Simple global models fail on complex data
\end{itemize}

\item Idea:
\begin{itemize}
\item Partition input space
\item Fit simple models locally
\end{itemize}

\item Implementation:
\begin{itemize}
\item Recursive binary splits
\item Optimize fit subject to complexity constraints
\end{itemize}
\end{itemize}

\subsection{Bayesian Regression Trees}

\begin{itemize}
\item Model:
\[
y(x) = g(x; T, M) + \epsilon, \quad \epsilon \sim \mathcal{N}(0,\sigma^2)
\]

\item Parameters:
\begin{itemize}
\item $T$: tree structure
\item $M$: leaf parameters
\end{itemize}

\item Prior factorization:
\[
\pi(T, M, \sigma^2) = \pi(M \mid T, \sigma^2)\pi(T \mid \sigma^2)\pi(\sigma^2)
\]
\end{itemize}

\subsection{Bayesian Additive Regression Trees (BART)}

\begin{itemize}
\item Model:
\[
y(x) = \sum_{j=1}^m g(x; T_j, M_j) + \epsilon
\]

\item Interpretation:
\begin{itemize}
\item Ensemble of weak learners
\item Each tree captures partial structure
\end{itemize}

\item Prior:
\[
\pi((T_1,M_1),\ldots,(T_m,M_m),\sigma^2)
=
\prod_{j=1}^m \pi(M_j \mid T_j,\sigma^2)\pi(T_j \mid \sigma^2)\pi(\sigma^2)
\]
\end{itemize}

\subsection{Regression Tree Representation}

\begin{itemize}
\item Tree defines mapping:
\[
g(x;T,M) \mapsto \mu(x)
\]

\item Leaf parameters:
\[
M = (\mu_1,\ldots,\mu_b)
\]

\item Possible leaf models:
\begin{itemize}
\item Constant
\item Linear
\item Gaussian process
\end{itemize}
\end{itemize}

\subsection{MCMC for Tree Models}

\begin{itemize}
\item Sampling scheme:
\begin{itemize}
\item Sample tree structure $T$ via MH
\item Sample leaf parameters $M$ via Gibbs
\item Sample variance $\sigma^2$ via Gibbs
\end{itemize}
\end{itemize}

\subsubsection*{Tree Proposals (MH)}

\begin{itemize}
\item Grow: split a terminal node
\item Prune: collapse a subtree
\item Change: modify split rule
\item Swap: exchange split rules
\end{itemize}

\begin{itemize}
\item Ensures reversibility of the Markov chain.
\end{itemize}

\subsection{Pros and Cons}

\subsubsection*{Advantages}
\begin{itemize}
\item Highly flexible function approximation
\item Handles mixed data types naturally
\item Scales via parallel MCMC
\end{itemize}

\subsubsection*{Limitations}
\begin{itemize}
\item Poor mixing (due to discrete tree moves)
\item Limited interpretability
\item Can underrepresent uncertainty
\end{itemize}

\section{Random Partition Models}

\subsection{Random Probability Measures and Clustering}

\begin{itemize}
\item In Bayesian nonparametrics, we place priors on random distributions:
\[
F \sim p(F), \quad F = \sum_{h=1}^{\infty} w_h \delta_{m_h}
\]

\item Many such priors (e.g., DP) are discrete $\Rightarrow$ induce ties.

\item Implied clustering:
\begin{itemize}
\item Draw $\theta_i \sim F$, i.i.d.
\item Positive probability that $\theta_i = \theta_{i'}$
\item This defines clusters via shared values
\end{itemize}

\item Partition:
\[
S_j = \{i : \theta_i = \theta_j^*\}, \quad j = 1,\dots,k
\]

\item Applications:
\begin{itemize}
\item Text: words grouped into topics
\item Biostatistics: patients grouped into subpopulations
\end{itemize}
\end{itemize}

\subsection{Inference on Partitions}

\begin{itemize}
\item If $F$ is not of direct interest:
\[
\rho = \{S_1,\dots,S_k\}
\]

\item Prefer specifying a prior directly on partitions:
\[
p(\rho)
\]
\end{itemize}

\subsection{Notation for Random Partitions}

\begin{itemize}
\item Units: $i \in S = \{1,\dots,n\}$

\item Partition:
\[
\rho_n = \{S_1,\dots,S_k\}, \quad S = \bigcup_{j=1}^k S_j
\]

\item Cluster sizes:
\[
n_j = |S_j|
\]

\item Allocation representation:
\[
s_i = j \iff i \in S_j
\]

\item Covariate-dependent partitions:
\[
p(\rho_n \mid x)
\]
\end{itemize}

\subsection{Sampling Model}

\begin{itemize}
\item Conditional on partition:
\[
p(y^n \mid \rho, \theta)
=
\prod_{j=1}^k \prod_{i \in S_j} p(y_i \mid \theta_j)
\]

\item Components:
\begin{itemize}
\item Cluster parameters $\theta_j$
\item Prior on $\theta_j$: often conjugate
\item Prior on $\rho$: defines clustering structure
\end{itemize}
\end{itemize}

\subsection{Product Partition Model (PPM)}

\begin{itemize}
\item Partition prior:
\[
p(\rho_n) \propto \prod_{j=1}^k c(S_j)
\]

\item $c(S_j)$: cohesion function measuring cluster strength

\item Examples:
\begin{itemize}
\item Dirichlet process:
\[
c(S_j) = \alpha (n_j - 1)!
\]

\item Contiguity (changepoint models):
\[
c(S_j) = p^{n_j - 1}(1-p)\mathbf{1}\{S_j \text{ contiguous}\}
\]
\end{itemize}

\item Conjugacy:
\begin{itemize}
\item Posterior remains PPM with updated cohesion
\end{itemize}
\end{itemize}

\subsection{Species Sampling Models (SSM)}

\begin{itemize}
\item Partition probability depends only on cluster sizes:
\[
p(\rho_n) = p(|S_1|,\dots,|S_k|)
\]

\item Exchangeable partition probability function (EPPF)

\item Predictive characterization (PPF):
\[
p(s_{i+1} = j \mid s_1,\dots,s_i)
\]

\item Must satisfy consistency conditions for exchangeability
\end{itemize}

\subsection{Polya Urn (Dirichlet Process)}

\begin{itemize}
\item Let $k_i$ be number of clusters among first $i$ observations

\item Predictive rule:
\[
p(s_{i+1} =
\begin{cases}
j & \text{existing cluster} \\
k_i + 1 & \text{new cluster}
\end{cases})
=
\begin{cases}
\frac{n_j}{M + i} \\
\frac{M}{M + i}
\end{cases}
\]

\item Equivalent to:
\begin{itemize}
\item DP prior on $G$
\item Special case of PPM and SSM
\end{itemize}
\end{itemize}

\subsection{Model-Based Clustering}

\begin{itemize}
\item Mixture model:
\[
p(y_i \mid k,\theta,\pi) = \sum_{j=1}^k \pi_j f_j(y_i \mid \theta_j)
\]

\item Latent assignment:
\[
p(y_i \mid s_i = j, \theta) = f_j(y_i \mid \theta_j), \quad \Pr(s_i=j)=\pi_j
\]

\item Induces partition via latent variables $s_i$
\end{itemize}

\subsection{Covariate-Dependent Clustering}

\subsubsection*{Motivation}

\begin{itemize}
\item Clustering should depend on covariates $x_i$
\item Similar $x_i$ $\Rightarrow$ more likely to co-cluster
\end{itemize}

\subsubsection*{Example 1: Sarcoma Trial}

\begin{itemize}
\item 12 cancer subtypes
\item Covariate: prognosis category
\item Goal: estimate success probabilities
\end{itemize}

Issues:
\begin{itemize}
\item Separate analysis: too little data
\item Full pooling: oversmoothing
\item Regression: insufficient flexibility
\end{itemize}

Solution:
\begin{itemize}
\item Cluster subtypes
\item Share strength within clusters
\item Let data determine clustering
\end{itemize}

\subsubsection*{Example 2: Patient Clustering}

\begin{itemize}
\item Cluster patients by response trajectories

\item Challenge:
\begin{itemize}
\item Want clustering informed by covariates
\item But also predictive for new observations
\end{itemize}
\end{itemize}

\subsection{Augmented Response Approach}

\begin{itemize}
\item Define augmented response:
\[
\tilde{y}_i = (y_i, x_i)
\]

\item Issues:
\begin{itemize}
\item Complex likelihood construction
\item Incorrect modeling of covariate influence
\end{itemize}

\item Better:
\begin{itemize}
\item Model partition directly via $p(\rho \mid x)$
\end{itemize}
\end{itemize}

\subsection{Mixtures with Covariate-Dependent Weights}

\begin{itemize}
\item Model:
\[
y_i \mid x_i \sim \sum_{j} \pi_j(x_i) f_j(y_i \mid \theta_j)
\]

\item Properties:
\begin{itemize}
\item Partition induced via latent $s_i$
\item Dependence enters through $\pi_j(x)$
\end{itemize}

\item Limitation:
\begin{itemize}
\item Fixed number of components (typically)
\end{itemize}
\end{itemize}

\subsection{Dependent Dirichlet Process (DDP)}

\begin{itemize}
\item Define predictor-indexed random measures:
\[
G_x = \sum_h \pi_h(x)\delta_{\theta_h}
\]

\item Features:
\begin{itemize}
\item Shared atoms $\theta_h$
\item Predictor-dependent weights $\pi_h(x)$
\end{itemize}

\item Induces covariate-dependent partitions
\end{itemize}

\subsection{DDP-like Constructions}

\subsubsection*{Probit Stick-Breaking Process}

\begin{itemize}
\item Weights:
\[
\pi_h(x) = \Phi(\eta_h(x)) \prod_{\ell<h} \{1 - \Phi(\eta_\ell(x))\}
\]

\item $\eta_h(x)$ controls dependence on covariates
\end{itemize}

\subsubsection*{Weighted Mixture of DPs}

\begin{itemize}
\item Combine local DPs:
\[
G_x(B) = \sum_{l=1}^L \frac{\gamma_l f(x,x_l^*)}{\sum_{l'} \gamma_{l'} f(x,x_{l'}^*)} G_l^*(B)
\]

\item Shared atoms, location-dependent weights
\end{itemize}

\subsection{Product Partition with Covariates}

\begin{itemize}
\item Modify cohesion:
\[
p(\rho) \propto \prod_{j=1}^k c(S_j) \exp\{t \cdot g(S_j)\}
\]

\item $g(S_j)$ measures similarity within cluster

\item Example:
\[
g(S_j) =
\frac{2}{|S_j|(|S_j|-1)}
\sum_{i<i'} \left(1 - \frac{d_{ii'}}{(1+\epsilon)d^*}\right)
\]

\item Encourages clustering of nearby points
\end{itemize}

\subsection{Summary}

\begin{itemize}
\item Random partitions provide a direct way to model clustering

\item Key frameworks:
\begin{itemize}
\item PPM: cohesion-based
\item SSM: size-based
\item DP: Polya urn
\item Model-based: latent mixtures
\end{itemize}

\item Extensions:
\begin{itemize}
\item Covariate-dependent partitions
\item Dependent random measures
\item Attraction-based clustering
\end{itemize}
\end{itemize}

\subsection{Item Allocation and Attraction-Based Clustering}

\begin{itemize}
\item Modify predictive probability function (PPF):
\[
p(s_{i+1} \mid s_1,\dots,s_i)
=
\begin{cases}
j & \propto \frac{1}{M+i} \, q(i+1, S_j) \\
k_i+1 & \propto \frac{M}{M+i}
\end{cases}
\]

\item Attraction function:
\[
q(i, S) = \sum_{\ell \in S} a_{i\ell}
\]

\item Define similarity weights:
\begin{itemize}
\item Distance $d_{ij}$ between items
\item $a_{ij} \propto [(1+\epsilon)d^* - d_{ij}]^T$
\item Normalize: $\sum_{j \neq i} a_{ij} = n-1$
\end{itemize}

\item Result:
\begin{itemize}
\item Encourages clustering of nearby points
\item Induces ergodic Markov chain (irreducible, aperiodic, recurrent)
\end{itemize}
\end{itemize}

\subsection{PPM with Covariates (PPMx)}

\begin{itemize}
\item Goal:
\[
p(\rho) \rightarrow p(\rho \mid x)
\]

\item Partition prior:
\[
p(\rho_n \mid x^n) \propto \prod_{j=1}^k g(x_j^*) \, c(S_j)
\]

\item $g(x_j^*)$: similarity function for cluster $S_j$

\item Construction:
\[
g(x_j^*) = \int \prod_{i \in S_j} q(x_i \mid \xi_j) q(\xi_j)\, d\xi_j
\]

\item Interpretation:
\begin{itemize}
\item Measures within-cluster covariate homogeneity
\item Penalizes heterogeneous clusters
\end{itemize}
\end{itemize}

\subsection{Properties of PPMx}

\begin{itemize}
\item Symmetry:
\[
p(\rho_n \mid x^*) \text{ invariant to permutations}
\]

\item Coherence across $n$:
\[
p(\rho_n \mid x^n)
=
\sum_{s_{n+1}} \int p(\rho_{n+1} \mid x^{n+1}) q(x_{n+1} \mid x^n)\, dx_{n+1}
\]

\item Requires:
\[
q(x_{n+1} \mid x^n) = \frac{g_{n+1}(x^{n+1})}{g_n(x^n)}
\]
\end{itemize}

\subsection{Choice of Similarity Models}

\begin{itemize}
\item Continuous covariates:
\[
q(x_i \mid \xi_j) = \mathcal{N}(\xi_j, V)
\]

\item Categorical:
\begin{itemize}
\item Multinomial likelihood
\item Dirichlet prior
\end{itemize}

\item Ordinal:
\begin{itemize}
\item Multinomial probit models
\end{itemize}

\item Count data:
\[
q(x_i \mid \xi_j) = \text{Poisson}(\xi_j), \quad q(\xi_j)=\text{Gamma}(a,b)
\]
\end{itemize}

\subsection{Posterior Inference under PPMx}

\begin{itemize}
\item Augmented model:
\[
q(y^n, x^n, \theta, \xi, \rho_n)
=
\prod_j \prod_{i \in S_j}
p(y_i \mid \theta_j, x_i)
q(x_i \mid \xi_j)
\, p(\theta_j)\, q(\xi_j)\, c(S_j)
\]

\item Posterior:
\[
p(\rho_n \mid x^n, y^n) \propto p(y^n, x^n \mid \rho_n)
\]

\item Computation:
\begin{itemize}
\item Reduces to standard clustering after augmentation
\item Straightforward Gibbs updates
\end{itemize}

\item Warning:
\begin{itemize}
\item Joint modeling of $(x,y)$ is not equivalent to PPMx
\end{itemize}
\end{itemize}

\subsection{Example: Mixed Effects Model}

\begin{itemize}
\item Data:
\begin{itemize}
\item $n=47$ patients
\item repeated measures $y_{ij}$
\item covariate $x_i$ (dose)
\end{itemize}

\item Model:
\begin{itemize}
\item Cluster patients
\item Share parameters within clusters
\end{itemize}

\item Sampling:
\[
p(y_i \mid \theta_j^*, \rho_n) = \mathcal{N}(\theta_j^*, \Sigma)
\]

\item Partition prior:
\[
p(\rho_n \mid x^n) \propto \prod_j c(S_j) g(x_j^*)
\]

\item Similarity:
\begin{itemize}
\item Multivariate normal model
\end{itemize}
\end{itemize}

\subsection{Covariate-Based vs Augmented Response}

\begin{itemize}
\item Augmented approach:
\begin{itemize}
\item Treat $(x,y)$ jointly
\item Clustering influenced by both
\end{itemize}

\item Covariate-based clustering:
\begin{itemize}
\item Partition depends only on $x$
\item Response modeled conditionally
\end{itemize}

\item Key difference:
\begin{itemize}
\item Augmented models distort similarity structure
\item PPMx preserves interpretability of clustering
\end{itemize}
\end{itemize}

\chapter{Latent Feature Models and the Indian Buffet Process (IBP)}

\section{Infinite Limit Construction}

\subsection{Key Idea: Used vs Unused Features}

\begin{itemize}
\item Partition columns into:
\begin{itemize}
\item Used: $m_k > 0$
\item Unused: $m_k = 0$
\end{itemize}

\item Define:
\[
K_+ = \text{number of used features}
\]

\item Reorder columns so:
\[
m_k > 0 \text{ for } k \le K_+, \quad m_k = 0 \text{ for } k > K_+
\]

\item Factorization:
\begin{itemize}
\item Product over unused columns repeats
\item Product over used columns carries structure
\end{itemize}
\end{itemize}

\subsection{Algebraic Simplification}

\begin{itemize}
\item Using Gamma identities:
\[
\frac{\Gamma(N+1+\alpha/K)}{\Gamma(1+\alpha/K)}
=
\prod_{j=1}^N \left(j + \frac{\alpha}{K}\right)
\]

\item And:
\[
\frac{\Gamma(m_k+\alpha/K)}{\Gamma(\alpha/K)}
=
\left(\frac{\alpha}{K}\right)\prod_{j=1}^{m_k-1}\left(j+\frac{\alpha}{K}\right)
\]
\end{itemize}

\subsection{Limiting Distribution ($K \to \infty$)}

\begin{itemize}
\item Harmonic number:
\[
H_N = \sum_{j=1}^N \frac{1}{j}
\]

\item Limit:
\[
P([Z])
=
\frac{\alpha^{K_+}}{\prod_h K_h!}
e^{-\alpha H_N}
\prod_{k=1}^{K_+}
\frac{(N-m_k)!(m_k-1)!}{N!}
\]

\item Interpretation:
\begin{itemize}
\item $\alpha^{K_+}$ rewards new features
\item $e^{-\alpha H_N}$ controls total growth
\item Exchangeable distribution
\end{itemize}
\end{itemize}

\section{Indian Buffet Process (IBP)}

\begin{itemize}
\item Infinite buffet interpretation:

\item Customer 1:
\[
K_1 \sim \text{Poisson}(\alpha)
\]

\item Customer $i$:
\begin{itemize}
\item Chooses existing feature $k$ with probability:
\[
\frac{m_k}{i}
\]

\item Draws new features:
\[
\text{Poisson}\left(\frac{\alpha}{i}\right)
\]
\end{itemize}

\item Result:
\begin{itemize}
\item Infinite features available
\item Finite number active per observation
\end{itemize}
\end{itemize}

\section{Posterior Inference}

\begin{itemize}
\item Finite model:
\[
P(z_{ik}=1 \mid Z_{-i,k}, X)
\propto
P(z_{ik}=1 \mid Z_{-i,k}) P(X \mid Z)
\]

\item Beta-Bernoulli marginal:
\[
P(z_{ik}=1 \mid Z_{-i,k})
=
\frac{m_{-i,k} + \alpha/K}{N + \alpha/K}
\]

\item Infinite limit:
\[
P(z_{ik}=1 \mid Z_{-i,k})
=
\frac{m_{-i,k}}{N}
\]

\item New features:
\[
\text{Poisson}(\alpha/N)
\]
\end{itemize}

\section{Properties of IBP}

\begin{itemize}
\item Exchangeable over observations

\item Row sums:
\[
\# \text{features per observation} \sim \text{Poisson}(\alpha)
\]

\item Total number of nonzeros:
\[
\sim \text{Poisson}(N\alpha)
\]

\item Number of active features:
\[
K_+ \sim \text{Poisson}(\alpha H_N)
\]
\end{itemize}

\section{Two-Parameter Extension}

\begin{itemize}
\item Extension introduces $\beta$:

\[
\pi_k \sim \text{Beta}\left(\frac{\alpha \beta}{K}, \beta\right)
\]

\item Predictive rule:
\begin{itemize}
\item Existing features:
\[
\frac{m_k}{\beta + i - 1}
\]

\item New features:
\[
\text{Poisson}\left(\frac{\alpha \beta}{\beta + i - 1}\right)
\]
\end{itemize}

\item Effect:
\begin{itemize}
\item $\beta$ controls sparsity and feature reuse
\end{itemize}
\end{itemize}

\section{Beta Process Connection}

\begin{itemize}
\item Beta process:
\[
B \sim \text{BP}(c, B_0)
\]

\item Generates weights:
\[
B = \sum p_i \delta_{\omega_i}
\]

\item Bernoulli process:
\[
X \sim \text{BeP}(B)
\]

\item Interpretation:
\begin{itemize}
\item Beta process generates feature probabilities
\item Bernoulli process generates binary matrix $Z$
\end{itemize}
\end{itemize}

\section{Stick-Breaking Construction for Beta Process}

\begin{itemize}
\item Generate:
\[
\mu_k \sim \text{Beta}(\alpha,1)
\]

\item Define weights:
\[
\pi_k = \prod_{j=1}^k \mu_j
\]

\item Differences from DP:
\begin{itemize}
\item Weights do NOT sum to 1
\item Features are independent
\end{itemize}
\end{itemize}

\section{Conjugacy}

\begin{itemize}
\item Model:
\[
B \sim \text{BP}(c, B_0), \quad X_i \sim \text{BeP}(B)
\]

\item Posterior:
\[
B \mid X_{1:n}
\sim
\text{BP}\left(c+n, \frac{c}{c+n}B_0 + \frac{1}{c+n}\sum_{i=1}^n X_i\right)
\]
\end{itemize}

\section{IBP from the Beta--Bernoulli Process}

\begin{itemize}
\item Let:
\[
B \sim \text{BP}(c, B_0), \quad X_i \mid B \sim \text{BeP}(B)
\]

\item Predictive distribution:
\[
X_{n+1} \mid X_{1:n}
\sim
\text{BeP}\left(
\frac{c}{c+n} B_0 + \frac{1}{c+n}\sum_{i=1}^n X_i
\right)
\]

\item Interpretation:
\begin{itemize}
\item Existing feature $j$ selected with probability:
\[
\frac{m_{n,j}}{c+n}
\]

\item New features:
\[
\text{Poisson}\left(\frac{c\gamma}{c+n}\right)
\]
\end{itemize}

\item This recovers the IBP predictive rule.
\end{itemize}

\section{Hierarchical Beta Process}

\begin{itemize}
\item Model:
\[
B \sim \text{BP}(c, B_0), \quad A_j \sim \text{BP}(c', B), \quad X_{ij} \sim \text{BeP}(A_j)
\]

\item Interpretation:
\begin{itemize}
\item Global feature pool: $B$
\item Group-specific features: $A_j$
\item Observations draw from group-level features
\end{itemize}

\item Enables:
\begin{itemize}
\item Sharing features across groups
\item Group-specific variability
\end{itemize}
\end{itemize}

\section{Completely Random Measures (CRM)}

\begin{itemize}
\item A random measure $\mu$ assigns random mass:
\[
\mu(A) \ge 0, \quad A \subset \Theta
\]

\item Completely random if:
\[
A_1 \cap A_2 = \emptyset \Rightarrow \mu(A_1) \perp \mu(A_2)
\]
\end{itemize}

\subsection{Lévy--Khintchine Representation}

\begin{itemize}
\item Any CRM can be written as:
\[
\mu = \sum_{k=1}^\infty w_k \delta_{\theta_k}
\]

\item $(w_k, \theta_k)$ arise from a Poisson point process with Lévy measure:
\[
\nu(dw, d\theta)
\]
\end{itemize}

\section{Poisson Point Process}

\begin{itemize}
\item Number of points in region $A$:
\[
N(A) \sim \text{Poisson}(\nu(A))
\]

\item Independence:
\[
A_1 \cap A_2 = \emptyset \Rightarrow N(A_1) \perp N(A_2)
\]

\item Conditional distribution:
\begin{itemize}
\item Given $N(A)=k$, points are i.i.d. in $A$
\end{itemize}
\end{itemize}

\section{Gamma Process}

\begin{itemize}
\item Definition:
\[
\mu \sim \text{GaP}(cH, b)
\]

\item Lévy measure:
\[
\nu(dw, d\theta) = c w^{-1} e^{-bw} dw \, H(d\theta)
\]

\item Key idea:
\begin{itemize}
\item Produces positive weights
\item Infinite number of atoms
\end{itemize}
\end{itemize}

\subsection{Connection to Dirichlet Process}

\begin{itemize}
\item Normalize:
\[
G = \frac{\mu}{\mu(\Theta)}
\]

\item Then:
\[
G \sim \text{DP}(\alpha, H)
\]

\item Interpretation:
\begin{itemize}
\item Gamma process → unnormalized measure
\item Normalization → probability measure
\end{itemize}
\end{itemize}

\section{Beta Process}

\begin{itemize}
\item Definition:
\[
B \sim \text{BP}(c, H)
\]

\item Lévy measure:
\[
\nu(dw, d\theta) = c w^{-1}(1-w)^{c-1} dw \, H(d\theta)
\]

\item Properties:
\begin{itemize}
\item $w_k \in (0,1)$
\item Interpreted as feature probabilities
\end{itemize}
\end{itemize}

\section{Beta Process $\rightarrow$ Bernoulli Process $\rightarrow$ IBP}

\begin{itemize}
\item Given:
\[
B = \sum_{k=1}^\infty w_k \delta_{\theta_k}
\]

\item Generate features:
\[
z_{ik} \mid B \sim \text{Bernoulli}(w_k)
\]

\item Feature vector:
\[
X_i = \sum_k z_{ik} \delta_{\theta_k}
\]

\item Marginalizing $B$:
\begin{itemize}
\item induces IBP over $Z$
\end{itemize}
\end{itemize}

\section{Interpretation Across Processes}

\begin{itemize}
\item Gamma process:
\begin{itemize}
\item weights on $(0,\infty)$
\item normalized → Dirichlet process
\end{itemize}

\item Beta process:
\begin{itemize}
\item weights on $(0,1)$
\item used as feature probabilities
\end{itemize}

\item Bernoulli process:
\begin{itemize}
\item generates binary features
\end{itemize}

\item IBP:
\begin{itemize}
\item marginal distribution over binary matrix $Z$
\end{itemize}
\end{itemize}

\section{Posterior Inference for IBP Models}

\begin{itemize}
\item Conditional update:
\[
P(z_{ik}=1 \mid Z_{-i,k}, X)
\propto
\frac{m_{-i,k}}{N} \cdot P(X \mid Z)
\]

\item New features:
\[
\text{Poisson}(\alpha/N)
\]

\item Structure:
\begin{itemize}
\item Columns independent
\item Rows exchangeable
\end{itemize}
\end{itemize}

\section{Example: Linear Gaussian Latent Feature Model}

\begin{itemize}
\item Data:
\[
x_i \in \mathbb{R}^d
\]

\item Model:
\[
x_i \sim \mathcal{N}(z_i A, \sigma^2 I)
\]

\item Where:
\begin{itemize}
\item $z_i$ binary feature vector
\item $A$ feature loading matrix
\end{itemize}

\item Matrix form:
\[
X = ZA + \epsilon
\]

\item Remains well-defined as $K \to \infty$
\end{itemize}

\section{Key Properties of IBP Models}

\begin{itemize}
\item Infinite-dimensional feature space
\item Finite features per observation
\item Exchangeable rows
\item Flexible latent representation
\end{itemize}


\end{document}